\documentclass[aps,pra,twocolumn,superscriptaddress,longbibliography]{revtex4-2}

\usepackage{amsmath,amssymb,amsthm,graphicx,xcolor,bm,braket,dcolumn,booktabs}
\usepackage{hyperref}

\newtheorem{theorem}{Theorem}
\newtheorem{conjecture}{Conjecture}
\newtheorem{observation}{Observation}

\begin{document}

\title{Toward PK-Free Molecular Hamiltonians:\\
Two-Center Sturmian Integrals and Coupled Composition}
\thanks{Paper 19 in the Geometric Vacuum series}

\author{J.~Loutey}
\affiliation{Independent Researcher, Kent, Washington}

\date{April 6, 2026}

%% ========================================================================
%% ABSTRACT
%% ========================================================================

\begin{abstract}
The composed natural geometry framework (Paper~17) constructs
molecular Hamiltonians by partitioning electrons into single-center
blocks and coupling them via Phillips--Kleinman (PK) pseudopotentials.
This produces structurally sparse qubit Hamiltonians---$O(Q^{2.5})$
Pauli scaling, 51--1{,}712$\times$ fewer terms than Gaussian
baselines (Paper~14)---but PK limits accuracy to 5.3--26\%
equilibrium bond length error.  PK cannot be improved: six
modification attempts have failed (v2.0.6--23), and the overcounting
diverges at higher angular momentum.

We report a negative result (Track~CB, v2.0.37): replacing PK with
explicit cross-block electron-electron integrals \emph{without}
corresponding cross-center nuclear attraction integrals produces an
unbalanced Hamiltonian (29\% error, worse than PK's 15\%).  The root
cause is identified: two-center one-body integrals
$\langle \psi_{nlm}^A | V_B | \psi_{n'l'm'}^A \rangle$, where an
orbital on center~$A$ feels the nuclear potential of center~$B$, are
required for energy balance but absent from the single-center
framework.

We observe that these missing integrals are precisely the
Shibuya--Wulfman integrals of the Coulomb Sturmian
literature~\cite{shibuya1965,avery2004,avery2012}, which are computed
via Fock's momentum-space projection onto $S^3$---the same conformal
map that underlies the GeoVac framework (Paper~7).  The mathematical
foundation is therefore already in place.  We review the convergence
properties of the Shibuya--Wulfman expansion, identify the Coulomb
resolution technique~\cite{hoggan2011} as a potentially faster
alternative, and outline the ``coupled composition'' architecture
that would result from incorporating two-center integrals.  If the
angular expansion converges in few terms at molecular bond lengths,
the coupled Hamiltonian would be PK-free, structurally sparse, and
significantly more accurate than the current composed framework.
This paper presents the balanced coupled architecture, documents its mathematical foundations, and reports computational benchmarks across three molecules (LiH, BeH$_2$, H$_2$O).

We report computational results from the first three phases.  The
convergence concern is resolved: the multipole expansion of the
cross-center potential terminates exactly at $L_{\max} = 2l_{\max}$
by Gaunt selection rules (33 one-body terms at $n_{\max}=2$).  The
balanced Hamiltonian has 878 Pauli terms (Conjecture~1 confirmed).
Four-electron FCI gives $R_{\rm eq}$ error of 7.0\% and
$D_e = 0.037$~Ha---the only bound configuration among all four
coupling schemes tested (Conjecture~2 not confirmed in strict sense,
but the comparison crosses FCI regimes).  At $n_{\max} = 3$ ($Q = 84$),
single-point FCI gives 0.20\% energy error (down from 1.8\%) with
analytical integrals (incomplete gamma functions, machine precision),
confirming convergence.  The $R_{\rm eq}$ drift is structural to the
block decomposition ($+0.053$~bohr per $n_{\max}$ step, $3\times$
smaller than PK's $+0.15$--$0.22$~bohr/$l_{\max}$), not numerical.
\end{abstract}

\maketitle

%% ========================================================================
\section{Introduction: The PK Problem}
\label{sec:intro}
%% ========================================================================

The composed natural geometry framework~\cite{loutey_paper17}
partitions electrons into groups---core, bond pairs, lone pairs---each
inhabiting its own natural coordinate system.  Inter-group coupling is
approximated by the Phillips--Kleinman (PK) pseudopotential, a
one-body operator that provides core-valence orthogonality and an
effective repulsive barrier.

PK is the framework's accuracy bottleneck.  At $l_{\max} = 2$ (the
optimal operating point), LiH achieves 5.3\% $R_{\rm eq}$ error;
BeH$_2$ gives 11.7\%; H$_2$O gives 26\%.  These errors are
structural, not numerical: the $l_{\max}$ divergence
(+0.15--0.22~bohr per additional angular channel) is confirmed to
originate in PK overcounting, not in the adiabatic approximation
(Track~BQ, v2.0.32).  Six attempts to modify PK have failed
(projector, spectral, self-consistent, $R$-dependent, $l$-dependent,
$l$-dependent on 2D solver), and the ab initio $R_{\rm ref}$
derivation failed because the relevant scale (${\sim}3$~bohr) is
molecular, not atomic (Track~AD, v2.0.18).

At the same time, PK is essential for the composed framework's
sparsity: it makes the electron repulsion integral (ERI) tensor
block-diagonal, yielding $O(Q^{2.5})$ Pauli scaling
(Paper~14)~\cite{loutey_paper14}.  Without PK, the block
decomposition has no coupling at all.

The question motivating this paper is: \emph{can PK be replaced by
explicit inter-block integrals that preserve structural sparsity while
eliminating the accuracy ceiling?}


%% ========================================================================
\section{The Negative Result: Coupled Composition Without $V_{ne}$}
\label{sec:negative}
%% ========================================================================

Track~CB (v2.0.37) tested the simplest version of this idea: remove
PK from the composed LiH Hamiltonian and add explicit cross-block
two-electron integrals $\langle pq | rs \rangle$ where $p, q$ belong
to one block and $r, s$ to another.

The cross-block ERIs were computed using existing infrastructure
(\texttt{cross\_block\_mp2.py}, originally built for Track~BX-3b).
Key census results for LiH at $n_{\max} = 2$, $Q = 30$:

\begin{table}[h]
\caption{Cross-block ERI census for LiH at $n_{\max} = 2$.}
\label{tab:census}
\begin{ruledtabular}
\begin{tabular}{lrr}
Metric & Within-block & Cross-block \\
\hline
Nonzero ERIs & 195 & 130 \\
Max $|\langle pq|rs\rangle|$ (Ha) & --- & 0.891 \\
Mean $|\langle pq|rs\rangle|$ (Ha) & --- & 0.142 \\
\end{tabular}
\end{ruledtabular}
\end{table}

The cross-block/within-block ERI ratio is 2/3, consistent with the
selection-rule counting bound of $\le 2\times$ from Paper~14,
Sec.~V.E.  Gaunt selection rules ($|\Delta l| \le 1$,
$|\Delta m| \le 1$) are preserved across block boundaries.

The resulting qubit Hamiltonian has 854 Pauli terms (2.56$\times$ the
composed value of 334) and 1-norm 85.7~Ha (2.30$\times$ the composed
value of 37.3~Ha).  However, the FCI energy tells the real story:

\begin{table}[h]
\caption{FCI accuracy for LiH at $n_{\max} = 2$ under different
coupling schemes.  Exact: $-8.071$~Ha.}
\label{tab:fci}
\begin{ruledtabular}
\begin{tabular}{lrr}
Configuration & $E_{\mathrm{FCI}}$ (Ha) & Error \\
\hline
No PK, no cross-block (decoupled) & $-7.195$ & 10.9\% \\
Composed (PK, no cross-block) & $-6.863$ & 15.0\% \\
Coupled (no PK, with cross-block $V_{ee}$) & $-5.734$ & 29.0\% \\
\end{tabular}
\end{ruledtabular}
\end{table}

The coupled Hamiltonian is \emph{worse} than both alternatives.  The
root cause is clear: cross-block ERIs add inter-group electron-electron
\emph{repulsion} without the balancing cross-center electron-nucleus
\emph{attraction}.  The Hamiltonian is energetically unbalanced.

\begin{observation}
\label{obs:decoupled}
The decoupled configuration (no PK, no cross-block interactions)
gives 10.9\% error---better than PK's 15.0\%.  This implies that
PK overcorrects at $n_{\max} = 2$: the blocks are more accurate
when left independent than when coupled via PK.
\end{observation}

This observation is consistent with the $l_{\max}$ divergence
diagnostic: PK's repulsive barrier systematically overcounts the
core-valence interaction energy, and the overcounting worsens with
basis size.


%% ========================================================================
\section{The Missing Piece: Two-Center Nuclear Attraction}
\label{sec:twocenter}
%% ========================================================================

The coupled Hamiltonian of Sec.~\ref{sec:negative} is incomplete.
A physically balanced molecular Hamiltonian requires, for each
electron:
\begin{enumerate}
\item Kinetic energy (within-block, already present).
\item Same-center nuclear attraction $V_A = -Z_A / r_A$ (within-block, already present).
\item Cross-center nuclear attraction $V_B = -Z_B / r_B$, where the
  electron on center~$A$ feels the potential of nucleus~$B$ (missing).
\item Electron-electron repulsion, both within-block (present) and
  cross-block (added in Track~CB).
\end{enumerate}

Item~3 requires two-center one-body integrals of the form
\begin{equation}
\label{eq:vne_cross}
I_{nlm, n'l'm'}^{AB} = \left\langle \psi_{nlm}^A \left|
  \frac{-Z_B}{|\mathbf{r} - \mathbf{R}_B|} \right|
  \psi_{n'l'm'}^A \right\rangle
\end{equation}
where $\psi_{nlm}^A$ is a hydrogenic orbital centered on nucleus~$A$,
and the potential is from nucleus~$B$ at position~$\mathbf{R}_B$.
These integrals are standard in the Slater-type orbital (STO)
literature but have never been computed in the GeoVac framework,
which has until now been strictly single-center at each level.


%% ========================================================================
\section{Connection to Coulomb Sturmian Theory}
\label{sec:sturmian}
%% ========================================================================

GeoVac's hydrogenic orbitals are Coulomb Sturmians with the
identification $k = Z/n$, where $k$ is the Sturmian exponent.
This is not merely a notational correspondence: the entire GeoVac
framework---the graph Laplacian on $S^3$ (Paper~7), the Fock
projection, the angular momentum selection rules---is built on the
same mathematical structure that the Avery group has developed for
molecular Coulomb Sturmian calculations since the
1990s~\cite{avery2000,avery2004,avery2012,avery2014}.

\subsection{Shibuya--Wulfman integrals}

When a Coulomb Sturmian centered on $A$ is expanded in the basis
centered on $B$, the expansion coefficients are the Shibuya--Wulfman
integrals~\cite{shibuya1965}:
\begin{equation}
\label{eq:sw}
S_{\tau', \tau} = \int d^3\mathbf{r}\; \chi_{\tau'}^*(\mathbf{r} -
\mathbf{R}_B)\, \frac{k}{r_A}\, \chi_\tau(\mathbf{r} - \mathbf{R}_A)
\end{equation}
where $\chi_\tau$ denotes a Coulomb Sturmian with quantum numbers
$\tau = (n, l, m)$, and the weighting $k/r_A$ arises from the
potential-weighted orthogonality of the Sturmian basis.

In Fock's momentum-space representation, Coulomb Sturmians map to
hyperspherical harmonics on $S^3$, and Shibuya--Wulfman integrals
become overlap integrals of translated hyperspherical harmonics.
The translation is parameterized by
\begin{equation}
S \equiv k \cdot R_{AB}
\end{equation}
where $R_{AB} = |\mathbf{R}_A - \mathbf{R}_B|$ is the internuclear
distance and $k$ is the common Sturmian exponent.

The nuclear attraction integral of Eq.~(\ref{eq:vne_cross}) is
closely related to the Shibuya--Wulfman integral: both involve a
Coulomb potential centered on a different nucleus acting on an orbital.
Koga and Matsuhashi~\cite{koga1987} derived sum rules for nuclear
attraction integrals over hydrogenic orbitals that express them as
finite sums over Clebsch--Gordan coefficients---the same angular
momentum algebra that generates Gaunt integrals in the GeoVac
framework.

\emph{Heterogeneous exponents.}  The standard Shibuya--Wulfman
formalism assumes a common Sturmian exponent~$k$ across all centers.
GeoVac's composed blocks use different effective charges
$Z_{\mathrm{eff}}$ per block, giving different exponents
$k_b = Z_{\mathrm{eff},b}/n$ for each block~$b$.  The generalized
overlap integrals between Sturmians of different exponents exist in
the literature; Koga and Matsuhashi~\cite{koga1987} give sum rules
for the common-exponent case, and the convergence properties for
heterogeneous exponents may differ from the common-$k$ case.  Any implementation must use these generalized
integrals, and the convergence diagnostic (Sec.~\ref{sec:path})
must be performed with the actual heterogeneous exponents of the
LiH block decomposition, not with a uniform~$k$.

\subsection{The convergence challenge}

The Shibuya--Wulfman expansion has a well-documented convergence
limitation.  When a Sturmian on center~$A$ is expanded in the
angular basis of center~$B$, the expansion involves hyperspherical
harmonics of increasing angular momentum~$L$.  The rate of convergence
is controlled by the parameter $S = k \cdot R_{AB}$.

For $S \ll 1$ (nuclei close together relative to the orbital scale),
convergence is rapid.  For $S \gtrsim 3$ (typical molecular bond
lengths), convergence is slow, requiring many angular terms.  The
Avery group explicitly notes this: the expansion ``is very slowly
convergent if $S$ is large''~\cite{avery2012}.

For LiH at equilibrium ($R_{AB} \approx 3$~bohr), with $k \approx 1$
for the valence orbitals, $S \approx 3$.  This places the system
squarely in the slow-convergence regime.  Each additional angular
momentum channel in the expansion adds orbitals and therefore Pauli
terms to the qubit Hamiltonian.  If $L_{\max} = 5$--10 is required
for convergence, the Pauli count could approach the full Level~4N
encoding (${\sim}3{,}288$ terms), defeating the purpose of the block
decomposition.

\subsection{Resolution: multipole expansion termination}
\label{sec:multipole_resolution}

The convergence concern of the preceding subsection applies to the
Shibuya--Wulfman \emph{basis} expansion: expanding an orbital
centered on~$A$ in the complete Sturmian basis of center~$B$, which
requires summing over all angular momenta~$L$.  The cross-center
nuclear attraction integral of Eq.~(\ref{eq:vne_cross}) uses a
different expansion: the standard multipole expansion of
$1/|\mathbf{r}-\mathbf{R}_B|$ in Legendre polynomials $P_L(\cos\theta)$.

For this expansion, the angular integral involves
\begin{equation}
\int Y_{l_1 m}^*(\hat{r})\, P_L(\cos\theta)\, Y_{l_2 m}(\hat{r})\, d\Omega\,,
\end{equation}
which is nonzero only when $|l_1 - l_2| \le L \le l_1 + l_2$ and
$l_1 + L + l_2$ is even.  For $n_{\max} = 2$ ($l_{\max} = 1$), the
maximum~$L$ needed is~2 ($p$--$p$ pairs); for $n_{\max} = 3$
($l_{\max} = 2$), $L_{\max} = 4$.  The expansion terminates exactly
at $L_{\max} = 2 l_{\max}$---no truncation error.

Furthermore, $m_1 = m_2$ is required: the cross-center potential is
diagonal in the magnetic quantum number when the internuclear axis is
along~$z$.  Combined with the Gaunt selection rules, this means the
cross-center $V_{ne}$ adds only 33 new one-body matrix elements to
$h_1$ for LiH at $n_{\max} = 2$.  These are essentially free in
computational cost compared to the two-body integrals.

The original grid-based radial integrals (trapezoid quadrature,
$O(1/n_{\mathrm{grid}})$ convergence) achieved 0.18\% accuracy at
$n_{\mathrm{grid}} = 4000$ for $Z = 1$ orbitals, and $9 \times 10^{-8}$
for $Z = 3$.  The radial split integrals (inner and outer regions
separated at $R_{AB}$) are integrals of polynomial times exponential:
each $R_{nl}(r)$ contributes a polynomial in $r$ times
$\exp(-Zr/n)$.  The product $R_{nl} R_{n'l'} r^p$ is a finite sum of
terms $c_j r^j \exp(-\alpha r)$ where $\alpha = Z(1/n + 1/n')$.
Each split integral evaluates exactly via the regularized incomplete
gamma function (\texttt{scipy.special.gammainc}/\texttt{gammaincc}).
This replaces the $O(1/n_{\mathrm{grid}})$ trapezoid quadrature with
machine-precision evaluation---an instance of the algebraicization
program (Paper~18, Sec.~XII).

\subsection{The Coulomb resolution alternative}

A potentially more efficient route was developed by Hoggan and
collaborators~\cite{hoggan2011}.  The ``Coulomb resolution'' technique
decomposes multi-center integrals over exponential-type orbitals into
sums of one-electron overlap-like products, each involving orbitals on
at most two centers.  Crucially, these two-center integrals are
separable in prolate spheroidal coordinates---the same coordinate
system used at Level~2 of the GeoVac hierarchy (Paper~11).

The Coulomb resolution sidesteps the slowly convergent angular
expansion of Shibuya--Wulfman by working directly in the two-center
coordinate system where the integrals are most naturally expressed.
GeoVac already possesses prolate spheroidal infrastructure from
Paper~11, including the Neumann expansion machinery for $V_{ee}$
(Paper~12).  The question is whether this infrastructure can be
extended to compute the one-body cross-center integrals of
Eq.~(\ref{eq:vne_cross}).


%% ========================================================================
\section{The Coupled Composition Architecture}
\label{sec:architecture}
%% ========================================================================

If two-center integrals can be computed efficiently, the resulting
``coupled composition'' Hamiltonian would have the following structure
for a molecule with $B$ blocks:

\begin{equation}
\label{eq:coupled}
H = \sum_{b=1}^{B} H_b^{(1)} + \sum_{b=1}^{B} V_{ee}^{(b)} +
\sum_{b < b'} V_{ee}^{(b,b')} + \sum_{b=1}^{B} \sum_{a \ne a_b}
V_{ne}^{(b,a)}
\end{equation}

where $H_b^{(1)}$ is the one-body Hamiltonian within block~$b$
(kinetic + same-center nuclear attraction), $V_{ee}^{(b)}$ is the
within-block ERI, $V_{ee}^{(b,b')}$ is the cross-block ERI, and
$V_{ne}^{(b,a)}$ is the cross-center nuclear attraction (electron in
block~$b$ feeling nucleus~$a$ from a different center).

Compared to the current composed Hamiltonian:
\begin{itemize}
\item PK is \emph{eliminated entirely}.  Core-valence orthogonality is
  enforced by the occupation number representation in second
  quantization, not by a pseudopotential.
\item Cross-block ERIs are included, subject to Gaunt selection rules.
  Track~CB established that these add ${\sim}2/3$ of the within-block
  ERI count (130 vs.~195 for LiH at $n_{\max} = 2$).
\item Cross-center nuclear attraction is included via two-center
  integrals.  The number of one-body terms is bounded by the number
  of orbital pairs times the number of off-center nuclei.
\end{itemize}

\begin{conjecture}
\label{conj:sparsity}
The Gaunt selection rules that enforce ERI sparsity within each block
also constrain the cross-block ERIs and the angular expansion of the
cross-center nuclear attraction.  If the two-center integral expansion
converges in $L_{\max} \le 3$ angular terms at molecular bond lengths,
the coupled Hamiltonian should have $\lesssim 1{,}500$ Pauli terms for
LiH at $n_{\max} = 2$---roughly $4.5\times$ the current composed
count of 334, but well below the full Level~4N count of 3{,}288 and
potentially below the Gaussian STO-3G count of 907 after the PK
one-body terms are removed.
\end{conjecture}

This conjecture is falsifiable by computing the two-center integrals
and performing the Jordan--Wigner transformation.

\begin{conjecture}
\label{conj:accuracy}
With a balanced Hamiltonian (all one-body and two-body terms present,
no PK), the FCI accuracy on the coupled Hamiltonian should be
significantly better than the composed framework's 5.3\% $R_{\rm eq}$
error, bounded below by the hydrogenic FCI accuracy floor
(${\sim}0.35\%$ for He at $n_{\max} = 5$, Paper~FCI-A).
\end{conjecture}


%% ========================================================================
\section{Relationship to Existing GeoVac Infrastructure}
\label{sec:infrastructure}
%% ========================================================================

The coupled composition program connects to multiple existing
components:

\emph{Paper~7 (Dimensionless Vacuum):} The $S^3$ conformal projection
that defines GeoVac's graph Laplacian is identical to Fock's
1935 momentum-space mapping~\cite{fock1935}, which is the foundation
of Shibuya--Wulfman theory.  The 18 symbolic proofs of Paper~7
validate the same mathematical structure that generates the
Shibuya--Wulfman integrals.

\emph{Paper~11 (Prolate Spheroidal):} The prolate spheroidal
coordinate system and Neumann expansion machinery are directly
applicable to the Coulomb resolution technique for two-center
integrals.

\emph{Paper~12 (Algebraic $V_{ee}$):} The Neumann expansion of
$1/r_{12}$ in prolate spheroidal coordinates demonstrates that
multi-center Coulomb integrals can be decomposed into algebraically
evaluable components via angular momentum recurrence.

\emph{Paper~14 (Qubit Encoding):} Section~V.E already bounds
cross-block ERI counts at $\le 2\times$ within-block; Track~CB's
empirical ratio of 2/3 is tighter.  The coupled Hamiltonian would
feed directly into the Jordan--Wigner encoding pipeline.

\emph{Paper~16 (Chemical Periodicity):} The $S_N$ representation
theory and atomic classifier determine the block decomposition.
Coupled composition uses the same block structure---it changes how
blocks are coupled, not how they are defined.

\emph{Paper~18 (Exchange Constants):} PK is classified as a
``calibration exchange constant''---a transcendental quantity that
enters when projecting between the graph and continuous coordinates.
If coupled composition eliminates PK, it removes one exchange constant
from the framework, consistent with the algebraicization program
(Sec.~XII of Paper~18).

\emph{Track~BX-3b (Cross-Block MP2):} The \texttt{cross\_block\_mp2.py}
module already computes cross-block ERIs for LiH.  This infrastructure
would be reused for the cross-block ERI component of the coupled
Hamiltonian.


%% ========================================================================
\section{Computational Results}
\label{sec:results}
%% ========================================================================

We report results from Phases~1--3 of the research path outlined in
Sec.~\ref{sec:path}.

\subsection{Quantum resource comparison}

Table~\ref{tab:resources} compares the qubit Hamiltonian metrics for
the four coupling schemes at LiH $n_{\max} = 2$, $Q = 30$.

\begin{table}[h]
\caption{Qubit Hamiltonian resources for LiH at $n_{\max} = 2$,
$Q = 30$.  All values from Jordan--Wigner encoding.  The Gaussian
STO-3G baseline ($Q = 12$) is included for reference.}
\label{tab:resources}
\begin{ruledtabular}
\begin{tabular}{lddd}
Configuration & \multicolumn{1}{c}{Pauli} & \multicolumn{1}{c}{1-norm (Ha)} & \multicolumn{1}{c}{QWC} \\
\hline
Composed (PK)      & 334  & 37.3 & 21  \\
Balanced (CD)      & 878  & 74.1 & 87  \\
Coupled (CB)       & 854  & 85.7 & 88  \\
Gaussian STO-3G    & 907  & 34.3 & 273 \\
\end{tabular}
\end{ruledtabular}
\end{table}

The balanced Hamiltonian has 878 Pauli terms, confirming
Conjecture~\ref{conj:sparsity}: the count is well below the
conjectured bound of ${\sim}1{,}500$.  The 1-norm of the balanced
Hamiltonian (74.1~Ha) is \emph{lower} than the coupled CB value
(85.7~Ha) because the attractive cross-center $V_{ne}$ partially
cancels the repulsive cross-block $V_{ee}$.  The QWC group count
(87) remains far below the Gaussian baseline (273).

\subsection{Cross-center $V_{ne}$ census}

The cross-center nuclear attraction integrals (Sec.~\ref{sec:multipole_resolution})
add 33 nonzero one-body matrix elements to $h_1$ for LiH at
$n_{\max} = 2$:

\begin{table}[h]
\caption{Cross-center $V_{ne}$ census for LiH at $n_{\max} = 2$,
$R = 3.015$~bohr.}
\label{tab:vne_census}
\begin{ruledtabular}
\begin{tabular}{lddd}
Block $\to$ nucleus & \multicolumn{1}{c}{Nonzero} & \multicolumn{1}{c}{Max (Ha)} & \multicolumn{1}{c}{Trace (Ha)} \\
\hline
Core ($Z\!=\!3$) $\to$ H  & 11 & 0.373 & -1.646 \\
Val-Li ($Z\!=\!1$) $\to$ H & 11 & 0.328 & -1.182 \\
H-side ($Z\!=\!1$) $\to$ Li & 11 & 0.985 & -3.544 \\
\hline
Total & 33 & --- & -6.372 \\
\end{tabular}
\end{ruledtabular}
\end{table}

The large trace from the H-side block feeling the $Z = 3$ lithium
nucleus ($-3.544$~Ha) is the dominant contribution, consistent with
the physical expectation that the hydrogen electron is most strongly
affected by the nearby heavier nucleus.

\subsection{Four-electron FCI}

Table~\ref{tab:fci_balanced} presents FCI results at $R = 3.015$~bohr
for all four coupling schemes, including the new balanced configuration.
Nuclear repulsion is $V_{NN}$ only (no PK contribution).

\begin{table}[h]
\caption{Four-electron FCI for LiH at $n_{\max} = 2$,
$R = 3.015$~bohr.  Exact: $-8.071$~Ha.}
\label{tab:fci_balanced}
\begin{ruledtabular}
\begin{tabular}{lddl}
Configuration & \multicolumn{1}{c}{$E_{\mathrm{FCI}}$ (Ha)} & \multicolumn{1}{c}{Error} & Bound? \\
\hline
Balanced (CD)        & -7.924 & 1.8\%  & YES \\
Decoupled            & -7.195 & 10.9\% & NO  \\
Composed (PK)        & -6.863 & 15.0\% & NO  \\
Coupled CB           & -5.734 & 29.0\% & NO  \\
Exact                & -8.071 & 0\%    & YES \\
\end{tabular}
\end{ruledtabular}
\end{table}

The balanced configuration is the only coupling scheme that produces a
bound molecule.  The PES scan gives $R_{\rm eq} = 3.226$~bohr (7.0\%
error vs.\ exact 3.015~bohr) and $D_e = 0.037$~Ha (exact: 0.092~Ha).
The variational bound is respected at all $R$-points.

Conjecture~\ref{conj:accuracy} predicted that the balanced Hamiltonian
would achieve significantly better accuracy than PK's 5.3\%
$R_{\rm eq}$ error.  The result (7.0\%) does not confirm this in strict
sense.  However, the comparison is cross-regime: the 5.3\% figure
uses $l$-dependent PK at $l_{\max} = 2$ with continuous PES optimization,
while the current result uses $n_{\max} = 2$ FCI on a qubit Hamiltonian
with grid-computed integrals (0.18\% radial accuracy for $Z = 1$).
The qualitative result---the balanced coupled configuration is the
\emph{only} bound configuration among all four coupling schemes---is the
primary finding of this phase.

\subsection{Basis convergence: $n_{\max} = 3$}

At $n_{\max} = 3$ ($Q = 84$, $l_{\max} = 2$), the balanced coupled
Hamiltonian has 19{,}959 Pauli terms, 448.9~Ha 1-norm, and 2{,}298 QWC
groups.  The Pauli scaling exponent (2-point fit from $Q = 30$ to
$Q = 84$) is 3.03; the 1-norm exponent is 1.75.  The
Pauli-to-composed ratio (2.53$\times$) is consistent with
$n_{\max} = 2$ (2.63$\times$), indicating that the cross-block
overhead is a constant factor, not a growing penalty.

The cross-center $V_{ne}$ multipole expansion terminates exactly at
$L_{\max} = 4$ for $l_{\max} = 2$ ($d$--$d$ pairs), confirmed by
comparison to $L_{\max} = 10$ (max difference: machine zero).  The
total cross-center $V_{ne}$ trace is $-12.54$~Ha (168 nonzero terms).

The original grid-based integrals (trapezoid quadrature,
$O(1/n_{\mathrm{grid}})$ convergence) systematically underestimated
cross-center nuclear attraction by 0.09--0.47\% for diffuse $Z = 1$
orbitals.  Replacing with analytical evaluation via incomplete gamma
functions eliminates this error entirely (machine precision).  This
algebraicization is an instance of the exchange constant program
(Paper~18, Sec.~VII): the grid-based $V_{ne}$ was a calibration
exchange constant that proved eliminable once the right algebraic
structure was identified.  The
corrected $n_{\max} = 3$ energy at $R = 3.015$ is $-8.055$~Ha (0.20\%
error), improved from grid-based $-8.048$ (0.28\%).  The $R_{\rm eq}$,
however, moved only from 3.287 (grid) to 3.280 (analytical), confirming
that the outward drift is structural to the block decomposition, not
numerical.

\begin{table}[h]
\caption{Balanced coupled LiH convergence with basis size.
Exact energy at $R = 3.015$~bohr: $-8.071$~Ha.  Grid: trapezoid
quadrature ($n_{\mathrm{grid}} = 4000$).  Analytical: incomplete gamma
functions (machine precision).}
\label{tab:convergence}
\begin{ruledtabular}
\begin{tabular}{ldddddd}
\multicolumn{1}{c}{Method} & \multicolumn{1}{c}{$n_{\max}$} & \multicolumn{1}{c}{$Q$} &
\multicolumn{1}{c}{Pauli} & \multicolumn{1}{c}{1-norm (Ha)} &
\multicolumn{1}{c}{$E(3.015)$ (Ha)} & \multicolumn{1}{c}{Error} \\
\hline
Grid       & 2  & 30 & 878    & 74.1  & -7.924 & 1.8\%  \\
Analytical & 2  & 30 & 878    & 74.1  & -7.929 & 1.7\%  \\
Grid       & 3  & 84 & 19959  & 448.9 & -8.048 & 0.28\% \\
Analytical & 3  & 84 & 19959  & 448.9 & -8.055 & 0.20\% \\
\end{tabular}
\end{ruledtabular}
\end{table}

Table~\ref{tab:req_comparison} shows the definitive $R_{\rm eq}$
comparison:

\begin{table}[h]
\caption{$R_{\rm eq}$ comparison: grid vs.\ analytical integrals.
Exact: 3.015~bohr.}
\label{tab:req_comparison}
\begin{ruledtabular}
\begin{tabular}{ldddd}
\multicolumn{1}{c}{Method} & \multicolumn{1}{c}{$n_{\max}$} &
\multicolumn{1}{c}{$R_{\rm eq}$ (bohr)} &
\multicolumn{1}{c}{$R_{\rm eq}$ err} &
\multicolumn{1}{c}{$E(3.015)$ err} \\
\hline
Grid       & 2 & 3.226 & 7.0\% & 1.8\%  \\
Analytical & 2 & 3.227 & 7.0\% & 1.7\%  \\
Grid       & 3 & 3.287 & 9.0\% & 0.28\% \\
Analytical & 3 & 3.280 & 8.8\% & 0.20\% \\
Exact      & - & 3.015 & 0\%   & 0\%    \\
\end{tabular}
\end{ruledtabular}
\end{table}

The $n_{\max} = 3$ analytical PES (5 R-points):

\begin{table}[h]
\caption{Balanced coupled LiH PES at $n_{\max} = 3$ (analytical integrals).}
\label{tab:pes_nmax3}
\begin{ruledtabular}
\begin{tabular}{dd}
\multicolumn{1}{c}{$R$ (bohr)} & \multicolumn{1}{c}{$E$ (Ha)} \\
\hline
2.9   & -8.0494 \\
3.015 & -8.0545 \\
3.1   & -8.0571 \\
3.3   & -8.0592 \\
3.5   & -8.0563 \\
\end{tabular}
\end{ruledtabular}
\end{table}

\subsection{Summary of conjecture status}

Conjecture~\ref{conj:sparsity} (Pauli count): \textbf{confirmed}.
The balanced Hamiltonian has 878 Pauli terms, well below the
conjectured ${\sim}1{,}500$ bound.

Conjecture~\ref{conj:accuracy} (accuracy): \textbf{MIXED}.  Analytical
integrals (machine precision via incomplete gamma functions) confirm that
the $R_{\rm eq}$ drift is structural to the block decomposition, not
numerical: grid error contributed only 0.007~bohr to $R_{\rm eq}$ at
$n_{\max} = 3$.  The drift is $+0.053$~bohr per $n_{\max}$ step
($3\times$ smaller than PK's $+0.15$--$0.22$~bohr/$l_{\max}$) but in
the same outward direction.  Energy converges excellently:
$1.8\% \to 0.20\%$ ($9\times$ improvement) from $n_{\max} = 2$ to 3.
The balanced approach is optimal for single-point quantum simulation at
fixed geometries (0.20\% energy error), but not for PES-based equilibrium
geometry determination.


%% ========================================================================
\section{Research Path}
\label{sec:path}
%% ========================================================================

The following steps are required to test the conjectures of
Sec.~\ref{sec:architecture}:

\begin{enumerate}
\item \textbf{Two-center integral implementation.}  \emph{COMPLETE.}
  The convergence concern from the Shibuya--Wulfman basis expansion
  (Sec.~\ref{sec:sturmian}) was resolved: the cross-center $V_{ne}$
  uses the standard multipole expansion, which terminates exactly at
  $L_{\max} = 2 l_{\max}$ by Gaunt selection rules
  (Sec.~\ref{sec:multipole_resolution}).  For LiH at $n_{\max} = 2$,
  only 33 one-body matrix elements are needed.  Original grid numerical accuracy:
  0.18\% at $n_{\mathrm{grid}} = 4000$ for $Z = 1$ orbitals,
  $9 \times 10^{-8}$ for $Z = 3$.  Subsequently replaced by analytical
  evaluation via incomplete gamma functions (machine precision, zero
  quadrature).  The decision gate ($L_{\max} > 5$)
  is moot---the expansion terminates at $L_{\max} = 2$.

\item \textbf{Balanced coupled Hamiltonian.}  \emph{COMPLETE.}
  The balanced Hamiltonian (Eq.~\ref{eq:coupled} with all four
  interaction types) gives 878 Pauli terms, 74.1~Ha 1-norm, and
  87 QWC groups (Table~\ref{tab:resources}).
  Conjecture~\ref{conj:sparsity} confirmed.

\item \textbf{FCI benchmark.}  \emph{COMPLETE.}
  Four-electron FCI gives $E = -7.924$~Ha (1.8\% error),
  $R_{\rm eq} = 3.226$~bohr (7.0\%), $D_e = 0.037$~Ha.
  The balanced configuration is the only bound coupling scheme
  (Table~\ref{tab:fci_balanced}).  Conjecture~\ref{conj:accuracy}
  not confirmed in strict sense (7.0\% vs.\ 5.3\% PK), but the
  comparison is cross-regime.

\item \textbf{Scaling study.}  \emph{COMPLETE.}  $n_{\max} = 3$
  FCI with analytical integrals (incomplete gamma functions, machine
  precision): $E(3.015) = -8.055$~Ha (0.20\% error, down from 1.8\%
  at $n_{\max} = 2$; see Table~\ref{tab:convergence}).
  Pauli scaling exponent 3.03, 1-norm exponent 1.75.
  $R_{\rm eq} = 3.280$~bohr (8.8\% error), confirming structural drift
  of $+0.053$~bohr per $n_{\max}$ step ($3\times$ smaller than PK's
  $+0.15$--$0.22$~bohr/$l_{\max}$).  Grid error contributes only
  0.007~bohr to $R_{\rm eq}$ (negligible).  Classification: MIXED---energy
  converges excellently, $R_{\rm eq}$ drifts structurally.
  Extension to BeH$_2$ and H$_2$O: COMPLETE (v2.0.41--42).
\end{enumerate}

The critical risk is Step~1: if the angular expansion requires
$L_{\max} > 5$ for convergence, the Pauli count will approach the
full Level~4N encoding and the approach loses its structural advantage.
The Coulomb resolution route (approach~b) is the fallback if direct
expansion is too slow.


%% ========================================================================
\subsection{Polyatomic census}
\label{sec:polyatomic}
%% ========================================================================

The balanced coupled framework was extended to BeH$_2$ (collinear,
v2.0.41) and H$_2$O (non-collinear, v2.0.42), completing the
polyatomic census for all three composed molecules.

\begin{table}[h]
\centering
\begin{tabular}{lccccccc}
\hline
System & Blocks & $Q$ & Composed & Balanced & Pauli & $\lambda_\text{comp}$ & $\lambda_\text{bal}$ \\
       &        &     & Pauli    & Pauli    & ratio & (Ha)                  & (Ha) \\
\hline
LiH    & 2 & 30 & 334   & 878   & 2.63$\times$ & 37.3  & 74.1  \\
BeH$_2$ & 3 & 50 & 556   & 2,652 & 4.77$\times$ & 354.9 & 304.7 \\
H$_2$O & 5 & 70 & 778   & 5,798 & 7.45$\times$ & 361$^\dagger$ / 28,053 & 1,509.3 \\
\hline
\end{tabular}
\caption{Balanced coupled resource census. $^\dagger$H$_2$O composed electronic-only (PK partitioned).}
\label{tab:balanced_census}
\end{table}

BeH$_2$ four-electron FCI gives 10.7\% energy error with the
balanced Hamiltonian, compared to 27.4\% with PK and 10.9\%
decoupled.  The balanced 1-norm (304.7~Ha) is \emph{lower} than the
composed 1-norm (354.9~Ha), confirming that PK adds more 1-norm than
cross-center $V_{ne}$ for this system.

H$_2$O required a non-collinear generalization: the O--H bond vectors
make a 104.5$^\circ$ angle, so the cross-center $V_{ne}$ matrices
cannot be computed in a shared $z$-axis frame.  The solution is Wigner
$D$-matrix rotation of the real spherical harmonic basis: compute
$V_{ne}$ in the $z$-frame aligned with each bond vector, then rotate
via block-diagonal $D \cdot V_z \cdot D^T$ where $D$ is the rotation
matrix for real spherical harmonics (identity for $l = 0$, Cartesian
rotation permuted to $m = (-1, 0, 1)$ ordering for $l = 1$).  FCI is
infeasible for H$_2$O (10 electrons, $Q = 70$).  The lone pair
$V_{ne}$ integrals are well-behaved: all traces are below 3~Ha per
nucleus, confirming that the multipole expansion does not produce
unphysical couplings for lone pair blocks.

The Pauli ratio (balanced/composed) grows with block count:
$2.63\times$ (2-block LiH) $\to$ $4.77\times$ (3-block BeH$_2$)
$\to$ $7.45\times$ (5-block H$_2$O), because each block pair adds
cross-block ERIs and cross-center $V_{ne}$ terms.  The growth is
consistent with an $O(B^2)$ scaling in the number of blocks $B$.

The 1-norm advantage is regime-dependent.  For VQE/NISQ applications
where Pauli count and QWC groups dominate circuit cost, the composed
architecture with PK remains cheaper.  For fault-tolerant QPE where
the 1-norm determines the number of Trotter steps, the balanced
framework eliminates PK entirely.  This is decisive for molecules
with large nuclear charges: H$_2$O's PK contributed 98.7\% of the
total composed 1-norm (28,053~Ha), and the balanced 1-norm
(1,509.3~Ha) is $18.6\times$ cheaper.  Even compared to the
electronic-only composed 1-norm (361~Ha, with PK classically
partitioned), the balanced 1-norm is $4.18\times$ higher---but the
balanced Hamiltonian does not require the classical 1-RDM post-processing
step that PK partitioning demands.


%% ========================================================================
\subsection{The W1c-residual orthogonality wall (second-row hydrides)}
\label{sec:w1c_residual}
%% ========================================================================

For $Z \geq 11$ second-row hydrides (NaH, MgH$_2$, HCl, ...), the
composed-coupled framework treats the heavy atom's [Ne] core as
``frozen'': the 10 core electrons are encoded only via the Hartree
screening profile $Z_{\rm eff}^B(\rho)$ on the cross-center
$V_{ne}$ multipole expansion, not as an explicit FCI block. This
choice makes the qubit Hamiltonian tractable (NaH at $n_{\max} = 2$
fits in $Q = 20$ qubits with 2 valence electrons), but it
introduces two structurally separable walls:

\paragraph{W1c (Hartree screening of cross-$V_{ne}$).}
Without screening, the bare cross-$V_{ne}$ uses $Z_B = 11$ as if the
[Ne] core were not there, leading to ${\sim}10\times$ over-attraction
of the H valence orbital. Phase~C--W1c
(\texttt{geovac/cross\_center\_screened\_vne.py}, May~2026) replaces
the bare Coulomb $-Z_B/\rho$ with the screened tail
$-Z_{\rm eff}^B(\rho)/\rho - \tilde\Phi^B(\rho)$ (Newton's
shell-theorem decomposition), driven by the
\texttt{neon\_core.FrozenCore} machinery. NaH PES at $n_{\max} = 2$:
the cross-$V_{ne}$ magnitude is reduced by 5.4--6.0$\times$ across
$R \in [1.5, 10]$~bohr, exactly matching the diagnostic prediction.

\paragraph{W1c-residual (W1b-class orthogonality).}
After W1c, NaH FCI energy descent is reduced from
$-6.24$~Ha (bare) to $-0.36$~Ha (screened). The residual
$0.36$~Ha overattraction at $R = 2.5$ bohr (still $4.8\times$
deeper than the experimental NaH $D_e \approx 0.075$~Ha, with no
internal PES minimum) is the named ``W1c-residual orthogonality
wall'': the H valence orbital is non-orthogonal to the [Ne]-core
orbitals on the heavy atom, and the framework does not enforce
Pauli exclusion of the H valence basis against the frozen core
beyond what Hartree screening captures.

\paragraph{The Phillips--Kleinman cross-center engineering closure
(negative result).}
The named engineering closure for the W1c-residual is the
cross-center analog of the same-center Phillips--Kleinman
projection (Sec.~\ref{sec:relationship}; Paper~17,
Eq.~(\ref{eq:Vpk}) for the same-center case). The implementation
(\texttt{geovac/phillips\_kleinman\_cross\_center.py}, May~2026)
adds the standard PK barrier
$\Delta H_{pq}^{\rm PK} = \sum_c (E_v - E_c) S_{pc} S_{cq}$ to the
H-block one-body Hamiltonian, with cross-center overlaps $S_{pc}$
computed by 2D Gauss--Legendre $\times$ radial-grid quadrature
in the H-orbital frame, real-SH azimuthal selection
($m_p^{\rm signed} = m_c^{\rm signed}$) enforced exactly,
and Clementi--Roetti HF eigenvalues $E_c$ for the [Ne]-core
orbitals (Na: $E_{1s} = -40.479$~Ha, $E_{2s} = -2.797$~Ha,
$E_{2p} = -1.518$~Ha).

The result is HONEST NEGATIVE: PK cross-center reduces the
W1c-residual descent from $0.357$~Ha to $0.305$~Ha
(14.6\% improvement) at standard absolute-PK parameters
($E_v = 0$, purely repulsive) but does NOT close the wall:
NaH PES is still monotonically descending with no internal
equilibrium minimum.

\begin{center}
\begin{tabular}{lcccc}
\hline
Configuration & $E_{\min}$~(Ha) & $R_{\min}$~(bohr) & Descent (Ha) & Bound? \\
\hline
bare ($V_{ne}$ unscreened, no PK)        & $-171.097$ & 2.5 & 6.244 & no \\
W1c (screened, no PK)                    & $-163.316$ & 2.5 & 0.357 & no \\
PK only (bare $V_{ne}$, with PK)         & $-170.949$ & 2.5 & 6.097 & no \\
\textbf{W1c+PK (screened+PK)}            & \textbf{-163.264} & \textbf{2.5} & \textbf{0.305} & \textbf{no} \\
\hline
\end{tabular}
\end{center}

The decomposition shows W1c and PK are necessary engineering
closures but not jointly sufficient: the residual ${\sim}290$~mHa
overattraction lives in mechanisms beyond W1c+W1b, with candidate
sources including (i) inner-core penetration physics that the
screened multipole expansion handles correctly in form but may
overestimate in magnitude near the unscreened nucleus,
(ii) the hydrogenic $Z_{\rm orb} = 1$ basis on the heavy-atom
valence side being qualitatively wrong (it is not the actual
Na 3s orbital but a generic $Z_{\rm eff} = 1$ hydrogenic function),
(iii) the cross-block ERI treatment between Na valence and H
valence at mismatched effective charges. A proper basis on the
heavy-atom valence side---e.g.\ radial Schr\"odinger eigenstates
of the FrozenCore $Z_{\rm eff}(r)$ potential, available in
\texttt{neon\_core.\_solve\_screened\_radial}---is the most
likely next engineering target.

The framework currently treats second-row chemistry-solver
binding as an open architectural problem. The walls have been
named and named-individually-closed-or-not at the engineering
level (W1c yes, W1c-residual / W1b PK no, basis still open),
but the joint closure does not yet produce a bound NaH PES.

\paragraph{The screened-Schr\"odinger valence basis (h1 diagonal only,
Track 3, May 2026 --- HONEST NEGATIVE on binding, DIAGNOSTIC POSITIVE
on the next target).}
The named follow-on after the PK negative was to ``replace the
$Z_{\rm orb} = 1$ basis on the Na center with radial Schr\"odinger
eigenstates of the FrozenCore $Z_{\rm eff}(r)$ potential.'' Track~3
(\texttt{geovac/screened\_valence\_basis.py}, 25~tests passing) implements
the diagonal-only version of this substitution: the framework's
$h_1[i,i] = -Z_{\rm orb}^2 / (2 \cdot \text{block}\_n^2)$ for the
heavy-atom-side valence orbitals is replaced by the actual eigenvalue
of \texttt{\_solve\_screened\_radial\_log} (s-states with
\texttt{allow\_l0=True}) or \texttt{\_solve\_screened\_radial}
($l \geq 1$). Wired into \texttt{build\_balanced\_hamiltonian} via
\texttt{screened\_valence\_basis: bool = False}; bit-exact backward
compat preserved (LiH 878-Pauli regression bit-identical at machine
precision; first-row blocks have \texttt{n\_val\_offset = 0} and are
auto-skipped).

The eigenvalue substitutions are physically meaningful:
Na 3s = $-0.170$~Ha (vs hydrogenic $-0.500$~Ha; +0.330 Ha correction),
Na 4s = $-0.067$~Ha (vs $-0.125$~Ha; +0.058 Ha),
Na 4p = $-0.051$~Ha (vs $-0.125$~Ha; +0.074 Ha each, three $m_l$ states).
The total trace shift on the NaH valence is $+0.611$~Ha --- comparable
to the W1c-residual descent depth (0.357 Ha).

Despite the large physical correction, the 8-configuration NaH PES
test (bare/W1c/PK/SV combinations at max\_n=2, Q=20, 2e FCI)
shows \emph{all 8 configurations remain monotonically descending with
no internal equilibrium minimum}. The descent ladder updates to:
bare $6.244$~Ha $\to$ +W1c $0.357$~Ha (17.5$\times$ reduction) $\to$
+W1c+PK $0.305$~Ha (1.17$\times$ further) $\to$ +W1c+PK+SV $0.313$~Ha
(0.97$\times$ --- slight worsening, within numerical noise).

\textbf{The critical structural finding} (the substantive result of
Track 3, even though binding is not recovered): the screened-valence
correction is \emph{R-independent at the diagonal $h_1$ level}.
Eigenvalues of the FrozenCore $Z_{\rm eff}(r)$ potential are properties
of the on-center potential, not of the bond. The correction matrix
is bit-identical at $R = 2.5$ and $R = 10$~bohr (max abs diff
$5.55 \times 10^{-17}$, machine precision). A constant diagonal
$h_1$ shift cannot affect the descent depth, which is by definition
an $R$-dependent energy difference. \emph{The W1c-residual wall
is structurally not a diagonal-h1 problem.}

\textbf{Diagnostic positive identification of the genuine target}
(\texttt{debug/w1c\_residual\_nah\_track3\_diag.py}): testing the
cross-$V_{ne}$ integral with hydrogenic-shape labels at the
\emph{physical} $n$ (block\_n=1, $l=0$ $\to$ $n_{\rm phys}=3$ hydrogenic 3s
instead of 1s) shows the framework over-attracts by $0.81$~Ha at
$R = 2.5$ vs the physical-$n$ shape, and by only $0.14$~Ha at
$R = 10$. The differential, $-0.674$~Ha, is approximately
$1.9\times$ the W1c-residual descent (0.357 Ha) --- a
wavefunction-shape substitution would over-correct from the
current attraction profile and likely produce a binding minimum.

The actual Na 3s screened wavefunction has mean radius
$\langle r \rangle = 4.47$~bohr, while the hydrogenic Z=1 1s has
$\langle r \rangle = 1.50$~bohr. The two have \emph{opposite} shape
biases: framework basis is too tight (1.50 bohr); the actual
chemistry has a much more diffuse 3s. Cross-center $V_{ne}$ scales
strongly with this orbital extent; the diagonal eigenvalue
substitution does not move the wavefunction shape.

\textbf{Verdict on the named engineering closure}: the SV diagonal
correction is necessary architectural infrastructure (the eigenvalue
machinery is correct and physically meaningful) but not sufficient
for binding recovery. The genuine engineering closure is
\emph{wavefunction-shape substitution in cross-$V_{ne}$ and
within-block ERIs} --- not eigenvalue substitution. Two architectural
options remain:
\begin{enumerate}
  \item \textbf{Physical-$n$ hydrogenic relabeling (1-week sprint)}:
    rewrite the orbital-state lists so that block\_n=1 with $l=0$ on
    Na uses physical $n=3$ hydrogenic shape (Slater integrals,
    $V_{ne}$ integrals, all evaluated at the physical principal
    quantum number) at $Z_{\rm orb}$ chosen to match the asymptotic
    Na valence Z. Cheap, uses existing hydrogenic Slater
    infrastructure, but introduces a minor convention break.
  \item \textbf{Strict-Schmidt orthogonalization (multi-week sprint)}:
    explicitly Schmidt-subtract the heavy-atom core orbitals from the
    valence basis before computing all matrix elements; mathematically
    definitive but requires re-implementing $h_1$, ERI, and cross-$V_{ne}$
    in the orthogonalized basis with renormalization.
\end{enumerate}
The diagnostic-shape test suggests Option~1 alone closes most of the
descent and is the natural next named track.

\paragraph{Modular propinquity reframing: the W1c residual as a Hilbert
$C^*$-bimodule deformation distance (Sprint M-Y, 2026-05-23).}
The modular propinquity literature
(Latr\'emoli\`ere~\cite{latremoliere2019modular,latremoliere2021dual})
provides a natural setting for the W1c-residual wall: the framework's
NaH valence sub-block at $n_{\max}=2$ is a Hilbert
$C^*$-bimodule $\mathcal{M}$ over $(\mathcal{A}_L, \mathcal{A}_R)$,
where $\mathcal{A}_L$ is the H-centered multiplication algebra (the
H-valence carrier) and $\mathcal{A}_R$ is the Na-centered algebra
with FrozenCore structure (the Na-valence carrier).  Candidate
``pin states'' for the bimodule element $\xi \in \mathcal{M}$ are
the framework defaults (hydrogenic $Z_\text{orb}=1$) versus
physically-screened alternatives (FrozenCore-screened Na 3s plus
hydrogenic H 1s).  The natural two-axis $L/R$ decomposition of the
bimodule deformation distance
\[
  d_{\mathcal{M}}^{LR}(\xi_a, \xi_b)
  = \bigl(d_L^2(\xi_a, \xi_b) + d_R^2(\xi_a, \xi_b)\bigr)^{1/2}
\]
where $d_{L,R}$ are integrated against Gaussian probes centered at
H ($d_L$) and Na ($d_R$) respectively, identifies the
\emph{right-action axis} (Na-side wavefunction shape) as the dominant
carrier of the W1c residual.  The structural prediction $d_R/d_L \gg 1$
is confirmed numerically by the Track $\alpha$-Diagnostic
(memo \texttt{debug/sprint\_alpha\_1\_diagnostic\_memo.md}):
$d_R/d_L = 3.23$ for the load-bearing (framework hydrogenic) vs.\
(physical Na 3s) pair at the canonical probe width $\sigma = 1$~bohr,
and $d_R/d_L \in [1.32, 3.58]$ across $\sigma \in [0.5, 3.0]$~bohr,
$d_R/d_L = 3.23$ stable across $R \in [3.0, 4.0]$~bohr, and bit-stable
to 4 digits under FrozenCore grid refinement.  The qualitative ordering
``$d_R \gg d_L$'' is robust under all consistency checks; the precise
factor is probe-width-dependent.

\paragraph{Alkali-hydride scaling refinement (Sprint $\alpha$-Diagnostic).}
The bimodule diagnostic was applied to three alkali-hydrides
(LiH, NaH, KH).  The original M-Y prediction was a uniform monotone
scaling $d_R(M) \sim r^\text{phys}_\text{valence}(M) - 1.5$~bohr,
expecting Li $<$ Na $<$ K $<$ Rb $<$ Cs.  The measured ordering is
\textbf{LiH $\ll$ NaH $\approx$ KH} (not strict monotone):
$d_R(\text{LiH}) = 0.42$, $d_R(\text{NaH}) = 0.72$,
$d_R(\text{KH}) = 0.64$.  The mechanism is the
\emph{radial-node count}:
Li 2s has 1 radial node, Na 3s has 2, K 4s has 3.  The K 4s
inner-shell oscillations align partially with the hydrogenic
$Z=1$ 4s oscillation structure in the bond region, producing
L\textsuperscript{2}-cancellation in the bond-region integrand
$|R^M_\text{hyd} - R^M_\text{phys}|^2 w_R(z)$ that NaH's smoother
(only 2 nodes) wavefunctions do not exhibit.  The refined uniform
scaling is therefore monotone-with-saturation, not strictly linear
in $|\Delta r|$.  The prediction for RbH (4 nodes) and CsH
(5 nodes) is that they continue the plateau, matching the framework's
empirical chemistry catalogue: LiH binds in production at
$n_{\max}=2$; NaH, KH, RbH, CsH do not.

\paragraph{Mixed Slater-$n$ essential for screened valence orbitals
(Sprint $\alpha$-Multi-zeta, 2026-05-23).}
The named follow-on to the bimodule diagnostic is to substitute
multi-zeta Slater expansions of the physically-screened Na 3s and 3p
wavefunctions for the framework's hydrogenic $Z_\text{orb}=1$ basis.
The fit (\texttt{geovac/multi\_zeta\_orbitals.py} Z=11 entry, Track
$\alpha$-Multi-zeta) returned an architectural finding distinct from
BBB93~\cite{bunge1993}/Clementi--Roetti~1974 conventions: the standard
single-orbital expansion convention of fixing
$n_\text{slater}^k = n_\text{orbital}$ for every primitive is
\emph{inadequate} for screened-Schr\"odinger eigenstates with
significant core penetration.  The physical Na 3s has $R(r=0.011)
\approx 3.06$~bohr$^{-3/2}$ (inner-shell core penetration) plus two
radial nodes plus a diffuse outer tail spanning $r \in [0.01, 15]$~bohr.
A pure $n=3$ basis with prefactor $r^2$ cannot reach the inner
amplitude without driving $\zeta$ to very large values, which then
suppresses the outer tail.  Empirically:
pure-$n=3$ at $K=3$ gave overlap 0.9992 but only 1 radial node
(node-count failure); mixed-$n=[1,1,2,3,3]$ at $K=5$ gave overlap
0.999999, $L^2$ error $1.4 \times 10^{-6}$, correct 2-node structure,
all $|c_k| < 1$.  \emph{Both zeta count and Slater $n$ must be flexible}
for a faithful screened valence basis.

\paragraph{The W1c residual has a third layer (Sprint $\alpha$-PES,
2026-05-23, NEGATIVE-AT-CURRENT-BASIS).}
The named target of M-Y Path A is to substitute the multi-zeta
physical Na 3s/3p basis into the production cross-$V_{ne}$ kernel via
the architectural extension of \texttt{geovac/shibuya\_wulfman.py}
(multi-zeta dispatch on \texttt{compute\_cross\_center\_vne}) and
\texttt{geovac/balanced\_coupled.py} (\texttt{multi\_zeta\_basis: bool
= False} kwarg).  The implementation is correct
(12 regression tests pass, bit-exact backward compat preserved); the
Step~1 algebraic kernel differential at the Na 3s on-site cross-$V_{ne}$
diagonal at NaH $R_\text{eq} = 3.566$~bohr is $-0.135$~Ha (more than
$2\times$ the gate threshold; sign-consistent with the dissociation
limit at $R = 10$~bohr where the differential drops to
$-0.025$~Ha, $18.7\%$ of the $R_\text{eq}$ value).
\emph{Step 2 FCI verdict: bit-zero}.  Multi-zeta substitution shifts
the NaH 2-electron FCI ground-state energy by
$|\Delta E| < 4 \times 10^{-13}$~Ha at every $R \in \{2.5, 3.0, 3.5,
4.0, 5.0, 10.0\}$~bohr.

The mechanism is a substantive new finding: \emph{at NaH $n_{\max}=2$
with bare cross-$V_{ne}$, the 2-electron FCI ground state lives
entirely on the H side, not the Na side}.  The H 1s sees the
un-screened Na $Z=11$ nucleus at $R = 3.5$~bohr with on-site
cross-$V_{ne}$ attraction of $-3.13$~Ha (matching the classical
point-charge limit $-Z_\text{Na}/R$ within $5\%$), compared to
the Na 3s self-attraction of $-0.28$~Ha.  The lowest 5 eigenvalues
of the $h_1$ matrix are entirely on H-side states; the Na 3s
eigenvalue sits at $i = 5$ in the eigenspectrum, \emph{above} the
lowest 5 bonding/antibonding orbitals on the H side.  The
2-electron FCI ground state puts both electrons in the lowest H-side
orbital with no Na 3s occupation, so the Na 3s diagonal/basis shift
has \emph{no effect on the FCI eigenvalue}.

\paragraph{Three-layer structure of the W1c residual (consolidated).}
After the $\alpha$ arc, the W1c-residual orthogonality wall on
second-row alkali hydrides decomposes into three structural layers:

\begin{itemize}
  \item \textbf{Layer 1 — H/Na orthogonality.}
    Addressed by
    \texttt{geovac/phillips\_kleinman\_cross\_center.py}.
    Verdict: NEGATIVE on binding (14.6\% improvement, PES still
    descending).  Architecture retained; module is clean.
  \item \textbf{Layer 2 — Na-side wavefunction shape.}
    Addressed by \texttt{geovac/multi\_zeta\_orbitals.py} Z=11 entry
    plus dispatch in \texttt{shibuya\_wulfman.py} and
    \texttt{balanced\_coupled.py}.
    Verdict: Step~1 algebraic differential POSITIVE
    ($-0.135$~Ha at NaH $R_\text{eq}$); Step~2 FCI verdict bit-zero.
    Architecture retained; the FCI invisibility is the substantive
    new finding.
  \item \textbf{Layer 3 — FCI variational state H-side localization
    at finite basis} (NEW).  At $n_{\max}=2$, the FCI ground state
    structurally localizes on the H side; on-site $h_1$ corrections
    and on-site basis substitutions are bit-zero on the FCI
    eigenvalue.  The residual lives in the variational state itself,
    not in operator-level matrix elements.
\end{itemize}

\paragraph{Three named follow-on directions.}
\begin{enumerate}
  \item \textbf{NaH $n_{\max}=3$ test} ($\sim 1$~week, highest priority).
    Bundled with the screened-path multi-zeta combination (currently
    flagged with \texttt{NotImplementedError}).  Falsifies or confirms
    Layer~3 as a basis-size artifact: if the larger basis lets the FCI
    ground state acquire Na-side weight, multi-zeta is load-bearing
    and the $\alpha$-PES verdict flips POSITIVE.
  \item \textbf{Cross-$V_{ne}$ kernel-shape substitution}
    ($\sim 1$-$2$~weeks, parallel).  Track~3's diagnostic (2026-05-09)
    identified this as the original named target with a substantial
    $-0.674$~Ha differential at $R = 2.5$~bohr.  The $\alpha$-PES Step~1
    finding ($-0.135$~Ha) is a smaller-effect cousin; the
    kernel-shape substitution is mechanically a refactor of
    \texttt{shibuya\_wulfman.py}'s \texttt{\_radial\_split\_integral}
    to use physical-$n$ orbital decomposition on the heavy-atom side
    instead of $n_\text{block}$ hydrogenic.
  \item \textbf{W1c + multi-zeta combination by removing
    NotImplementedError} ($\sim 2$-$3$ days, bundled with item~1).
    Architectural unification of the bare and screened paths.
\end{enumerate}

\paragraph{F1 arc:\ unified W1c $\times$ multi-zeta architecture and the
structural success / energetic failure split (Sprints F1-P1+P2, W1c $\times$
M-Z bridge, and F1 max\_n=3, 2026-05-23).}
The F1 arc completed the multi-zeta integration at the architectural
level and tested the unified architecture at two basis sizes
(NaH max\_n=2 and max\_n=3). The arc closed with a CLEAN NEGATIVE on
binding and a substantive new structural finding that sharpens the
W1c-residual wall mechanism.

\textbf{F1-P1+P2 (unified architecture at max\_n=2):} The defensive
\texttt{NotImplementedError} blocking the combined W1c $\times$ multi-zeta
dispatch (in \texttt{geovac/balanced\_coupled.py} line 748) was removed
cleanly: the bare-Coulomb and screening-correction sub-integrals compose
at the integrand level, so the unified screened+multi-zeta path is a
${\sim}25$-line addition to \texttt{geovac/cross\_center\_screened\_vne.py}
with a \texttt{multi\_zeta\_basis: Optional[Dict] = None} kwarg. Three
findings at NaH max\_n=2:\ (a) multi-zeta is bit-zero on the FCI in the
bare cross-$V_{ne}$ case (reproducing the $\alpha$-PES Step~2 finding) but
\emph{fully load-bearing when combined with W1c} (differential $+0.056$~Ha
at $R = 3.5$~bohr scaling to $+0.208$~Ha at $R = 2.0$);\ (b) the original
Layer-3 framing was revised --- at W1c the FCI natural occupations are
$[1.0, 1.0]$ (open-shell-singlet-like with Na 3s occupation ${\sim}0.98$),
NOT H-dominant as inferred from the $\alpha$ arc;\ (c) the combined
architecture reduces descent depth from W1c-alone $0.898$~Ha to
W1c+mz $0.690$~Ha ($23\%$ reduction) but PES remains monotonically
descending. Verdict: PARTIAL-CLOSURE-AT-MAX\_N=2.

\textbf{W1c $\times$ M-Z bridge sprint (three-bucket hypothesis):} The bridge
sprint applied the modular propinquity M-Z partition (cross-shift vs
endomorphism; cf.\ Sub-sprint M-Z, 2026-05-23) to the W1c three-layer
hierarchy and proposed a three-bucket refinement: cross-shift
(handled) / basis-closable cross-shift (handled at sufficient basis
size) / endomorphism (external input). Sub-layer~3 (orbital-pair
flexibility) was classified as basis-closable cross-shift: the bonding
combination $[\text{H 1s} \pm \text{Na 3s}]_\pm$ is structurally a
bimodule cross-shift (left+right linear combination), conditionally
handled at sufficient basis size. Three falsifiable predictions for
NaH max\_n=3 were frozen into
\texttt{debug/data/sprint\_w1c\_mz\_partition\_predictions.json} before
the test ran:\ \textbf{P1} (internal PES minimum at $R_\text{eq} \in
[3.0, 4.5]$~bohr), \textbf{P2} (binding $D_e \in [0.0375, 0.150]$~Ha
within $2\times$ experimental $D_e \approx 0.075$~Ha), \textbf{P3}
(multi-zeta differential $\in [0.02, 0.30]$~Ha at $R_\text{eq}$).

\textbf{F1 max\_n=3 (CLEAN NEGATIVE).} The combined architecture at
$Q = 56$ ran the three predictions:

\begin{center}
\begin{tabular}{lll}
\hline
Prediction & Actual & Verdict \\
\hline
P1: internal min $R_\text{eq} \in [3.0, 4.5]$ bohr & $R_\text{min} = 2.0$ bohr (smallest tested) & \textbf{FAIL} \\
P2: $D_e \in [0.0375, 0.150]$ Ha & $D_e^\text{PES} = +0.7097$ Ha ($10\times$ exp.) & \textbf{FAIL} (cond.) \\
P3: $|E_\text{W1c} - E_\text{W1c+mz}| \in [0.02, 0.30]$ Ha & $|0.2066|$ Ha at $R = 2.0$ & \textbf{PASS} \\
\hline
\end{tabular}
\end{center}

P1 fails on the exact falsification criterion ($R_\text{min}$ at the
smallest tested $R$); the PES is monotonically descending across
$R \in \{2.0, 2.5, 3.0, 3.5, 4.0, 5.0, 7.0, 10.0\}$~bohr. P2 fails
(conditional) — the spurious-binding descent of $0.7097$~Ha is $\sim 10\times$
the experimental binding energy. P3 passes — multi-zeta endomorphism
content is preserved across basis sizes (consistent with max\_n=2's
$0.21$~Ha at the same $R$), confirming the endomorphism classification
of multi-zeta itself does not depend on basis size. Aggregate: 1 of 3
predictions pass; CLEAN NEGATIVE per the bridge sprint's gate.

\textbf{The substantive new structural finding (structural success /
energetic failure split).} The headline content the prediction test did
not anticipate: at max\_n=3 under combined W1c+mz, the dominant natural
orbital IS a true bonding combination. At $R = 3.5$:
\[
  \phi_\text{dom}
    \;=\; -0.698 \cdot \phi_\text{Na, 3s}
       \;-\; 0.687 \cdot \phi_\text{H, 1s}
       \;+\; \text{small admixtures}
\]
with Na/H amplitude$^2$ split $0.50/0.50$. The 2nd natural orbital is
the antibonding combination
$\phi_\text{2nd} = -0.698 \cdot \phi_\text{Na, 3s} + 0.687 \cdot \phi_\text{H, 1s}$.
At max\_n=2 the dominant NO was 100\% H-localized at every architecture;
at max\_n=3 the basis enlargement DID provide the orbital-mixing
flexibility predicted by the bridge sprint. The Na/H amplitude split is
stable at ${\sim}50/50$ across $R \in [2, 5]$~bohr --- the bonding
combination is not a localized feature of $R = 3.5$; it persists across
the PES.

\emph{But the constructed bonding orbital has higher energy than the
separated configuration at every tested $R$.} The PES is monotonically
descending; natural occupations remain $[1.0, 1.0]$ (open-shell
singlet of bonding $+$ antibonding singly-occupied), NOT $[2.0, 0.0]$
(closed-shell bond where the bonding pair is doubly occupied). The
framework can BUILD the right orbital but cannot energetically
PREFER it. The two halves of ``binding'' --- orbital constructibility
and energetic preference --- split cleanly at this resolution.

\textbf{Sub-layer 3 reclassification.} The bridge sprint's hypothesis
that sub-layer~3 is BIMODULE CROSS-SHIFT (basis-closable) is
\emph{partially confirmed and partially refuted}. The ``structural''
part is correct (the right linear combination IS structurally a bimodule
cross-shift, and the basis enlargement DOES make it constructible). The
``framework handles it'' part is too strong: the framework constructs
the orbital but cannot energetically prefer it. \textbf{Sub-layer~3
reclassifies to MODULE ENDOMORPHISM (basis-irreducible at tested
scales)}, joining sub-layer~2 (multi-zeta basis-shape) in the
endomorphism bucket. The proposed three-bucket M-Z partition is
FALSIFIED; the original two-bucket partition (cross-shift / endomorphism)
stands.

\textbf{F2 cross-$V_{ne}$ kernel-shape diagnostic (KERNEL-NOT-IT plus
architectural-absence finding).} Following the F1 max\_n=3 CLEAN NEGATIVE,
the next named target was the cross-$V_{ne}$ kernel-shape substitution on
the H partner side (Track~3's P2). Per the diagnostic-before-engineering
rule, F2 ran the Step~1 algebraic diagnostic before committing to a
1-2~week production extension: five matrix elements computed via converged
3D numerical quadrature vs the framework's default multipole expansion at
$L_\text{max}=4$. Largest converged differential: $2.0\times 10^{-5}$~Ha
on the H 1s diagonal cross-$V_{ne}$ contribution, which is $2.7\times 10^{-4}
= 0.027\%$ of the NaH bond-energy scale (0.075~Ha) --- six to ten orders of
magnitude better than what would be needed to close the W1c-residual wall.
\textbf{Verdict: KERNEL-NOT-IT.} The multipole expansion is bit-faithful where
it lives; further precision work within the existing block-diagonal
architecture cannot close the wall.

\emph{The substantive new content the gate did not anticipate}: the
framework's h1 in \texttt{composed\_qubit.build\_composed\_hamiltonian} is
\textbf{strictly block-diagonal}. There is NO cross-block off-diagonal h1
matrix element $\langle\psi_a^A \mid V_{ne}(\text{any nucleus}) \mid
\psi_b^B\rangle$ for orbitals on different centers $A \neq B$. The
framework constructs h1 as a strict block-diagonal sum of atom-like
sub-block Hamiltonians, with no slot for cross-block one-body coupling.
The bonding/antibonding splitting from h1 alone in the
\{$\ket{\text{Na 3s}}, \ket{\text{H 1s}}$\} basis is therefore identically
zero. The F1 max\_n=3 bonding orbital (constructible from natural-orbital
diagonalization of the 1-RDM) was emerging entirely from 2-body ERI
coupling in the FCI, not from h1 cross-coupling --- because the latter
does not exist in the architecture. Sub-layer~3 reclassifies (third time,
getting sharper): from F1's ``module endomorphism'' to F2's
\textbf{``architectural absence''} --- a missing matrix slot in the
existing block-diagonal architecture, distinct from a missing calibration
input or a wrong-value matrix element. New sub-wall name: \textbf{W1d}
(cross-block h1 architectural extension), structurally orthogonal to W1c
(frozen-core screening), W1a, W1b, W2a, W2b, W3.

\textbf{F3 cross-block h1 architectural extension (W1d CLOSED at FCI level
+ W1e NEWLY NAMED).} F3 implements the missing matrix slot: new module
\texttt{geovac/cross\_block\_h1.py} ($\sim 437$~lines, s-s only for first
pass, axial geometry only, 18 new tests). The cross-block off-diagonal
h1 matrix elements
$\langle\psi_a^A \mid T + \sum_C (-Z_C/|r-R_C|) \mid \psi_b^B\rangle$
for orbitals on different centers $A \neq B$ are computed via 2D axial
Gauss-Legendre quadrature with closed-form $\phi$ integral (s-orbital
geometry), auto-rotation of the frame so the two centers are collinear,
and kinetic via the Hellmann-Feynman identity
$T\psi_B = (E_B + Z_B/r_B)\psi_B$. Production code: \texttt{build\_balanced\_hamiltonian}
gains 5 \texttt{cross\_block\_h1} kwargs; \texttt{build\_composed\_hamiltonian}
gains 6 kwargs; bit-exact backward compat at default \texttt{cross\_block\_h1=False}
preserved (146 baseline tests + 18 new pass, zero regression).

\emph{Step~1 algebraic diagnostic at NaH $R_\text{eq}=3.566$~bohr:} for
hydrogenic Na~3s, total h1[Na 3s, H 1s] $= +0.320$~Ha, bonding/antibonding
splitting (lower-energy generalized eigenvalue) $3.46$~Ha. For multi-zeta
Na~3s: total $= -1.370$~Ha (sign flip; the multi-zeta Na 3s extends further
into the bond region with mean radius 4.5~bohr vs hydrogenic 1.5~bohr),
splitting $4.20$~Ha. Both satisfy the strong-bonding decision-gate
($50$~mHa threshold) by 1-2 orders of magnitude.

\emph{Step~3 FCI verdict at NaH max\_n=2 (Q=20, 2-electron singlet):}

\begin{center}
\begin{tabular}{lrrrl}
\hline
Architecture & $E(R=3.5)$ Ha & $E(R=10)$ Ha & $D_e$ (2-pt) Ha & Naturals (R=3.5) \\
\hline
bare & $-169.204$ & $-164.672$ & $+4.532$ & $[1.9933, 0.0067]$ (H-only bonding) \\
W1c+mz & $-163.115$ & $-162.779$ & $+0.336$ & $[1.0000, 1.0000]$ (separated, no NO bonding signature) \\
\textbf{W1c+mz+xblockh1} & $\mathbf{-167.767}$ & $\mathbf{-163.985}$ & $\mathbf{+3.782}$ & $\mathbf{[1.9991, 0.0007]}$ \textbf{(bonding, 50/50 Na/H)} \\
\hline
\end{tabular}
\end{center}

\textbf{The natural-orbital signature transformation is the F3 headline
structural finding.} W1c+mz alone has TWO singly-occupied separated
orbitals (naturals exactly $[1.0000, 1.0000]$). W1c+mz+xblockh1 has ONE
doubly-occupied bonding orbital (naturals $[1.9991, 0.0007]$). This is a
\textbf{four-orders-of-magnitude change} in the dominant natural-orbital
occupation signature, driven entirely by the cross-block h1 architectural
extension. The h1 lowest eigenvector character flipped concurrently from
Na/H $= 0.00/1.00$ (H-localized) at W1c+mz alone with eigvalue $-0.802$~Ha,
to Na/H $= 0.51/0.49$ (bonding combination) at W1c+mz+xblockh1 with
eigvalue $-3.161$~Ha --- $\sim 2.4$~Ha lower than the previously-lowest
antibonding-Na-localized state. \textbf{W1d (the architectural absence
named by F2) is CLOSED at the FCI level.}

\emph{Step~4 mini-PES with W1c+mz+xblockh1 (8 R-points across $[2, 10]$~bohr):}
PES is monotonically descending with $R_\text{min} = 2.0$~bohr (smallest
tested) and well depth $D_e^\text{F3} = +4.37$~Ha --- $58\times$
experimental NaH $D_e \approx 0.075$~Ha. The dominant NO is bonding (50/50
Na/H) at every R; the cross-block h1 construction is robust across the
full PES.

\emph{F3 verdict: PARTIAL-CLOSURE.} Bonding-orbital construction CLOSED;
inner-region overattraction OPENS. \textbf{New sub-layer W1e (inner-region
overattraction)}: cross-block h1 introduces a bonding mechanism into the
framework whose eigenvalue is lower than either separated atomic eigenvalue
and whose matrix elements grow as orbitals overlap (intensifies at smaller
R). The 2-electron FCI on NaH max\_n=2 has NO explicit frozen-core
electrons (the [Ne] core on Na is treated as a screening potential via
W1c, not as occupied orbitals); therefore no Pauli-exclusion mechanism
prevents the bonding pair from collapsing toward small R. This is the
structural analog of the missing-Pauli-repulsion mechanism that standard
Phillips-Kleinman pseudopotentials supply. The framework's existing PK
cross-center barrier operates on the partner-side valence orbital diagonal,
NOT on the bonding-combination orbital that cross-block h1 newly
constructs --- the natural F4 follow-on is the bonding-orbital PK
extension.

\textbf{Revised W1c wall mechanism (five-sub-layer hierarchy after the
F1$\to$F2$\to$F3 arc).}
\begin{itemize}\setlength\itemsep{0pt}
\item \textbf{W1c-cross-screening (CLOSED, Phase~C, 2026-05-08).}
  Frozen-core $Z_\text{eff}(r)$ reduces cross-$V_{ne}$ by $5$--$6\times$
  on second-row systems. Bimodule cross-shift handled at the operator
  level via the screened-multipole structure.
\item \textbf{W1c-multi-zeta-basis (CLOSED at architectural level, F1-P1+P2,
  2026-05-23).} Physical screened Na 3s replaces hydrogenic
  $Z_\text{orb}=1$ basis. Chemistry-side module endomorphism, mitigable via
  FrozenCore + atomic-physics fits; bit-zero without cross-shift activator
  (bare cross-$V_{ne}$), load-bearing with W1c on ($23\%$ additional descent
  reduction at max\_n=2).
\item \textbf{W1d-cross-block-h1 (CLOSED at FCI level, F3, 2026-05-23).}
  Off-diagonal h1 between orbitals on different centers, structurally
  absent from \texttt{composed\_qubit} prior to F3. Adding it via
  \texttt{geovac/cross\_block\_h1.py} transforms naturals
  $[1.0, 1.0] \to [1.9991, 0.0007]$ at NaH max\_n=2 (four orders of
  magnitude in dominant occupation signature), constructing the bonding
  combination at the FCI ground state. Architectural-absence category
  (F2's sharpening of sub-layer~3): a missing matrix slot in the existing
  block-diagonal architecture, structurally distinct from kernel-error and
  from calibration-data endomorphism, framework-internally closable via
  architectural extension.
\item \textbf{W1e-inner-region-overattraction (NEWLY NAMED, F3, 2026-05-23;
  F4 target).} Missing Pauli repulsion between the constructed bonding
  pair and the [Ne] frozen-core orbitals. Cross-block h1 lowers the bonding
  eigenvalue monotonically with decreasing R; no opposing repulsion in the
  2-electron architecture. Net well depth $D_e^\text{F3} = +4.37$~Ha vs
  experimental $0.075$~Ha. Natural F4 closure: extend
  \texttt{geovac/phillips\_kleinman\_cross\_center.py} (currently
  partner-side valence diagonal) to project the F3-constructed bonding
  combination onto the [Ne] frozen-core orbitals and apply Phillips-Kleinman
  repulsion at the bonding-orbital level.
\item \textbf{(FALSIFIED hypothesis: three-bucket M-Z refinement.)} The
  bridge sprint's basis-closable cross-shift sub-class evaporates at
  F1 max\_n=3: the bonding combination IS structurally a bimodule
  cross-shift and IS basis-closable at max\_n=3 in the orbital-construction
  sense, but the framework cannot energetically prefer it from within its
  bimodule machinery at the tested scales. Two-bucket M-Z partition stands
  with architectural-absence as a structurally distinct third
  sub-category WITHIN the partition.
\end{itemize}

\textbf{Three named F4 follow-on candidates.}
\begin{enumerate}\setlength\itemsep{0pt}
\item \textbf{F4 — Bonding-orbital Phillips-Kleinman extension (TOP PRIORITY,
  ${\sim}1$~week).} Extend
  \texttt{geovac/phillips\_kleinman\_cross\_center.py} to act on the
  cross-block-h1-constructed bonding orbital (rather than on the
  partner-side valence diagonal). After cross-block h1 is added in F3's
  phase~3.5, diagonalize h1 to identify the dominant bonding combination;
  compute its overlap with the [Ne] frozen-core orbitals via the
  F3-style 2D axial quadrature; add a PK projector
  $\Delta H^\text{PK,bond} = \sum_c (E_v - E_c)\cdot\mathcal{P}_\text{bond}
  S_{vc}\bra{c}$ at the operator level. Decision gate: internal minimum
  at NaH max\_n=2 with $R_\text{eq}$ within 1~bohr of 3.566 and well
  depth in $[0.0375, 0.150]$~Ha closes the W1c-residual wall to within
  $2\times$ experimental binding energy. Extends existing PK infrastructure
  rather than requiring new architecture.
\item \textbf{F3-extend: $l>0$ cross-block h1 + non-axial geometry
  ($\sim 2$~weeks).} Current F3 cross-block h1 is s-s only; mixed angular
  momentum requires spherical harmonic decomposition at both centers with
  their own quantization axes, bipolar harmonic algebra or full 3D
  Lebedev quadrature, and non-collinear geometry handling. Necessary to
  generalize F3's architectural advance from NaH/MgH$_2$/HCl alkali
  hydrides to polyatomic frozen-core systems (H$_2$S, PH$_3$, ...).
\item \textbf{DO NOT pursue further max\_n enlargement on NaH OR within-sub-block
  kernel-precision work.} F1 max\_n=3 + F2 + F3 between them closed the
  basis-size and kernel-precision directions: at max\_n=3 the structural
  gap (bonding orbital constructibility) is closed by basis enlargement;
  at max\_n=2 with cross-block h1 the structural gap is closed by
  architectural extension. The remaining wall (W1e) is NOT basis-size
  limited and NOT kernel-precision limited; it is a Pauli-repulsion gap at
  the bonding-orbital level, addressed only by F4 architectural extension.
\end{enumerate}

\paragraph{Source documentation.}
Track $\alpha$-Diagnostic:
\texttt{debug/sprint\_alpha\_1\_diagnostic\_memo.md}.
Track $\alpha$-Multi-zeta:
\texttt{debug/sprint\_alpha\_2\_multizeta\_memo.md}.
Track $\alpha$-PES:
\texttt{debug/sprint\_alpha\_3\_pes\_test\_memo.md}.
Track $\beta$ synthesis (modular + $\alpha$ arc):
\texttt{debug/sprint\_modular\_alpha\_arc\_synthesis\_memo.md}.
F1-P1+P2 unified architecture:
\texttt{debug/sprint\_f1\_p1p2\_combined\_test\_memo.md}.
W1c $\times$ M-Z partition bridge:
\texttt{debug/sprint\_w1c\_mz\_partition\_analysis\_memo.md}.
F1 max\_n=3 CLEAN NEGATIVE:
\texttt{debug/sprint\_f1\_maxn3\_predictions\_test\_memo.md}.
Sprint $\beta$-neg synthesis (F1-maturity, superseded by full-arc memo):
\texttt{debug/sprint\_beta\_neg\_synthesis\_memo.md}.
\textbf{F2 cross-$V_{ne}$ kernel-shape diagnostic (KERNEL-NOT-IT +
architectural-absence finding, W1d named):}
\texttt{debug/sprint\_f2\_cross\_vne\_kernel\_memo.md}.
\textbf{F3 cross-block h1 architectural extension (W1d CLOSED + W1e
newly named):}
\texttt{debug/sprint\_f3\_cross\_block\_h1\_memo.md}.
\textbf{Comprehensive full-arc synthesis (F1$\to$F2$\to$F3 at F3 maturity):}
\texttt{debug/sprint\_w1c\_full\_arc\_synthesis\_memo.md}.

\subsubsection{Refinement of the W1e mechanism after F4--F6}
\label{sec:w1e_refinement_f4_f6}

The F3 PARTIAL-CLOSURE verdict named W1e (inner-region overattraction)
as the natural F4 target with the working hypothesis ``missing
Pauli repulsion against the [Ne] frozen core,'' with three candidate
closure mechanisms ranked by priority: (1) bonding-orbital
Phillips--Kleinman extension (F4), (2) explicit frozen-core electrons
in FCI (F5 priority), (3) valence basis enlargement (F6 priority).
Three same-day diagnostic sprints (F4, F5, F6) tested all three
candidates and refined the W1e mechanism into a structurally distinct
multi-sub-mechanism reading. \textbf{None of the three candidates
closes W1e on its own; together they account for at most
${\sim}10$--$25\%$ of the wall, leaving ${\sim}65$--$90\%$ as
genuinely-deep correlation requiring architectural lifts at the
multi-week+ scale.}

\textbf{F4 --- bonding-orbital PK extension RULED OUT
(rank-1 PK saturation $\to 43\%$ closure ceiling).} The Step~1
algebraic diagnostic at NaH $R_\text{eq}=3.566$~bohr computes the
bonding-orbital--core overlap from the full F3 h1 diagonalization
($\max |S| = 17.5\%$ on the Na 2s core, well above the $5\%$
strong-gate threshold) and the corresponding Phillips--Kleinman
barrier at strict absolute reference $E_v = 0$ as $+0.194$~Ha
(22$\times$ smaller than the 4.37~Ha wall depth). The
diagnostic-before-engineering FCI sensitivity test injects a
rank-1 PK shift $\Delta H = c \cdot |\text{bond}\rangle\langle
\text{bond}|$ on the F3 h1 at varying $c$: at the Step~1-predicted
$c = +0.19$~Ha the FCI well depth changes by $0.05\%$
(4.3735~Ha $\to$ 4.3712~Ha), and even at $c = 2.0$~Ha (10$\times$
the predicted barrier) closure saturates at $43\%$ (residual
${\sim}2.5$~Ha, still 33$\times$ experimental). The rank-1
PK projection on the bonding orbital is structurally near-orthogonal
to the 2-electron FCI's dominant correlation channels at small $R$
--- the FCI ground state puts significant weight on configurations
OTHER than the bonding-pair doubly-occupied state, and a
rank-1 perturbation on the bonding orbital does not lift those
configurations. \textbf{W1e is therefore NOT a single-particle
Pauli-orthogonality wall} (the working hypothesis F3 named); it is
a multi-determinant FCI correlation effect that rank-1 PK cannot
suppress. F4 STOP-at-Step-1 documented in
\texttt{debug/sprint\_f4\_bonding\_pk\_memo.md}.

\textbf{F5 --- explicit-core Hartree treatment RULED OUT
($25.7\%$ closure ceiling at 2-point level).} The Step~1 algebraic
diagnostic computes the cross-block 2-body Coulomb correction
$\sum_c (J - K)_{\text{bonding},c} = +1.12$~Ha (Hartree $J = +1.13$~Ha,
Slater--Condon $K = +0.008$~Ha, $K/J = 0.72\%$) between the F3
bonding orbital and the five Na [Ne] core orbitals at NaH
$R_\text{eq} = 3.566$~bohr using Clementi--Raimondi exponents. The
correction is the right SIGN (repulsive, would lift the inner-region
collapse) but only $25.7\%$ of the 4.37~Ha wall depth. Linear
extrapolation of the predicted Step~3 outcome gives residual
$D_e^\text{F5} \approx +3.25$~Ha = 43$\times$ experimental NaH
$D_e \approx 0.075$~Ha, well outside the prespecified
$2\times$ closure window $[0.0375, 0.150]$~Ha. \textbf{The F4
precedent applies: F5 saturation at higher Hartree strengths would
probably reach $30$--$50\%$ closure, structurally insufficient.}
F5 STOP-at-Step-1 documented in
\texttt{debug/sprint\_f5\_explicit\_core\_memo.md}.
W1e reclassified from ``FCI correlation channels addressable by
explicit cores'' to \textbf{``deep correlation effect beyond
Hartree-level core-bonding J -- K interaction.''}

\textbf{F6 --- valence basis enlargement to max\_n=4 RULED OUT
($10.2\%$ PES closure ceiling; 2-point gate overestimates by 2.5$\times$).}
F6 ran the full F3 stack (W1c-screening + multi-zeta-basis +
cross-block-h1) at NaH max\_n=4 ($Q = 120$, $M = 60$, FCI dim 3600;
${\sim}5$~min/single-point) across 5 R-points $\{2.0, 3.0, 3.566,
5.0, 10.0\}$~bohr. PES well depth at $R_\text{min}=2.0$:
$D_e^\text{PES, F6} = +3.929$~Ha, \textbf{$10.2\%$ closure}. The
2-point gate (R=3.5 vs R=10) used in F3 §3 suggests $26.1\%$ closure
at max\_n=4 but this is \textbf{ARTIFACTUAL} --- the PES scan reveals
the inner-region collapse continues from R=3.5 down to R=2.0 (energy
drops by another 0.696~Ha that the 2-point gate misses). The bonding
signature is structurally preserved at every $R$ (dominant NO
$\in [1.9971, 1.9992]$; Na/H amplitude $\approx 50/50$). Methodological
finding for future quick-gate testing in this regime: \textbf{the
2-point gate overestimates PES-derived closure by 2.5$\times$ on
this class of monotonically-descending PES; future sprints should
use PES well depth at $R_\text{min}$ as the load-bearing closure
metric}, not 2-point energy differentials at fixed reference geometry.
F6 PARTIAL\_IMPROVEMENT documented in
\texttt{debug/sprint\_f6\_maxn4\_nah\_memo.md}.

\textbf{Refined six-sub-layer W1c hierarchy after F4--F6.}
W1e decomposes into three structurally INDEPENDENT sub-sub-mechanisms,
each with a measured closure ceiling within sprint-scale reach:
\begin{itemize}\setlength\itemsep{0pt}
\item \textbf{W1e-basis-truncation:} ${\sim}10.2\%$ closure ceiling
  at NaH max\_n=4 PES level (F6, measured); diminishing returns
  expected at max\_n=5+ based on the 10\% level at the 6$\times$
  basis-size increase from max\_n=2.
\item \textbf{W1e-Hartree-pressure:} ${\sim}25.7\%$ closure ceiling
  at max\_n=2 2-point level (F5, predicted from cross-block 2-body
  Coulomb $J - K$ Hartree mean-field); 2-pt-vs-PES caveat from F6
  suggests PES-level closure may be lower than the 2-pt prediction.
\item \textbf{W1e-deep-correlation-residual:} ${\sim}65$--$90\%$ of
  the wall remains beyond mean-field + single-determinant +
  basis-enlargement architecture. Closure candidates: Schmidt
  orthogonalization at the basis-construction level (${\sim}3$--$4$
  weeks), explicit treatment of [Ne] core as fully correlated
  (multi-month, requires occupancy-restriction FCI infrastructure
  not present in framework), or acceptance that the framework's
  structural-skeleton scope (CLAUDE.md §1.7) does not autonomously
  extend to binding energetics at the second-row regime within
  current compute budgets.
\end{itemize}

\textbf{Methodological caution flagged for future sprints
(2-point-vs-PES discrepancy).} F1, F2, F3 all used a 2-point gate
at $R = 3.5$~bohr vs $R = 10$~bohr as the quick-gate test for $D_e$.
F6's PES scan reveals this gate overestimates wall depth by
${\sim}2.5\times$ for monotonically-descending PES in the W1c-residual
regime, because the inner-region collapse extends below $R = 3.5$
down to $R_\text{min} = 2.0$ with significant additional descent
(${\sim}0.7$~Ha at max\_n=4) that the 2-point gate misses. \textbf{Future
quick-gate tests in this regime should use a 3-point or 5-point
mini-PES across $R \in [2, 10]$~bohr} rather than a 2-point
differential at fixed reference geometry, to avoid mis-attributing
closure fractions. F3's reported $D_e^\text{F3} = +4.37$~Ha is itself
a 2-pt value (R=3.5/R=10); the corresponding F3 PES well depth at
$R_\text{min}=2.0$ is likely larger (untested but expected by analogy
to F6).

\textbf{Net W1e closure verdict after F4--F6: deep correlation that
resists basis-truncation, single-particle PK, and mean-field
Hartree-class closure at scales reachable within sprint-scale compute.}
The chemistry arc's forward direction shifts from ``find the right
closure mechanism'' to ``characterize the framework's
structural-skeleton scope limits empirically for second-row hydride
binding.'' Consistent with the broader CLAUDE.md §1.7
multi-focal-composition wall taxonomy finding (the framework's
structural-skeleton scope does not autonomously deliver binding
energetics at this precision).

\paragraph{Source documentation (F4--F6 refinement).}
F4 bonding-orbital PK STOP-at-Step-1:
\texttt{debug/sprint\_f4\_bonding\_pk\_memo.md}.
F5 explicit-core Hartree STOP-at-Step-1:
\texttt{debug/sprint\_f5\_explicit\_core\_memo.md}.
F6 max\_n=4 PARTIAL\_IMPROVEMENT:
\texttt{debug/sprint\_f6\_maxn4\_nah\_memo.md}.
Sprint $\beta$-update F4-F6-APPENDED-TO-F3-BASELINE synthesis:
\texttt{debug/sprint\_beta\_update\_f4f6\_memo.md}.

\subsubsection{Closure of the W1e sprint-scale candidate space and cross-domain pattern (Schmidt + core correlation, 2026-05-23)}
\label{sec:w1e_schmidt_core_correlation}

The two named multi-week+ candidates from \S\ref{sec:w1e_refinement_f4_f6}
(``Schmidt orthogonalization at basis-construction level'' and
``fully correlated [Ne] cores'') were tested by Day-1 diagnostics,
executing the diagnostic-before-engineering rule at sprint-cycle scale.
\textbf{Both ruled out as W1e closure mechanisms.}

\textbf{Schmidt orthogonalization at basis-construction level
(Day-1 algebraic diagnostic, STOP).} The H 1s basis function on the
H center has total projection mass
$\sum_c |S_c|^2 = 1.00 \times 10^{-3}$ on the Na [Ne] core at
NaH $R = 3.5$~bohr --- only $0.1\%$ of H 1s lives in [Ne]-core
space. Three nonzero overlaps (Na 1s, 2s, 2p$_0$) survive the
axial $m$-selection rule; Na 2s dominates
($|S|^2 = 8.9 \times 10^{-4}$, $|S| \approx 3.0\%$). The
$17.5\%$ max-overlap F4 found
(\S\ref{sec:w1e_refinement_f4_f6}) was specifically on the
F3-constructed \emph{bonding orbital}, which mixes H 1s with
Na 3s; the H 1s basis function alone is structurally far outside
the [Ne] core's significant amplitude. Total diagonal-element
differential: $\Delta_\text{total} = +8.6$~mHa (converging from
below across a 4-step Cartesian-grid sweep $80^3 \to 140^3$), which
is $\sim 12\times$ below the 100~mHa prespecified GO threshold
and $\sim 3$--$4\times$ below the 30~mHa BORDERLINE threshold;
it would close at most $\sim 0.2\%$ of the 4.37~Ha W1e wall depth.
\textbf{The W1e wall is not a basis-level non-orthogonality wall.}
Day-1 driver and memo:
\texttt{debug/sprint\_w1e\_schmidt\_diagnostic\_memo.md},
\texttt{debug/data/sprint\_w1e\_schmidt\_diagnostic.json}.

\textbf{Fully correlated [Ne] cores (Day-1 literature
order-of-magnitude estimate, DEFER).} The binding-relevant
[Ne] core-correlation differential at NaH $R_\text{eq}$ converges
from three independent literature anchors to $\sim 0.001$--$0.003$~Ha
(central), $\sim 0.007$~Ha (upper bound): Iron-Oren-Martin 2003
qualitative ranking of Na subvalence correlation strength;
Sylvetsky-Martin 2018 W4-17 diatomic core-valence band; atomic
Na [Ne]-core correlation ($\sim 0.36$~Ha) times an estimated
active fraction at NaH bond geometry ($0.5$--$2\%$). All three
sit $14$--$1000\times$ below the prespecified $0.1$~Ha DEFER
threshold. \textbf{Core correlation is a Layer-2 sub-percent
effect for Na chemistry; it cannot close a $\sim 3.85$~Ha
structural wall (scales mismatch by $\sim 10^3$).} Day-1 memo:
\texttt{debug/sprint\_w1e\_core\_correlation\_estimate\_memo.md},
\texttt{debug/data/sprint\_w1e\_core\_correlation\_estimate.json}.

\textbf{Net W1e closure verdict: sprint-scale candidate space
exhausted.} Five independent sprint-scale W1e closure mechanisms
now ruled out: F4 single-particle Phillips--Kleinman ($43\%$
saturation ceiling), F5 mean-field Hartree $J - K$ ($25.7\%$
predicted ceiling), F6 basis enlargement at $n_{\max} = 4$
($10.2\%$ PES closure), basis-level Schmidt orthogonality
($0.2\%$ closure), [Ne] core correlation
($\sim 10^{-3}$--$10^{-2}$ of wall). The W1e wall is deep
multi-determinant correlation at the heavy-atom valence/core
boundary in the second-row regime; closure candidates beyond
sprint scale are multi-month architectural commitments
(self-consistent $Z_\text{eff}$ depending on cross-center
electronic density; fully-correlated all-electron treatment
incompatible with the two-electron active-space architecture;
fundamentally different basis families) or acceptance of the
framework's structural-skeleton scope as the autonomous boundary
for second-row binding energetics.

\textbf{Cross-domain pattern crystallization.} The W1e wall is the
\textbf{sixth structurally-independent instance of the
multi-focal-composition wall pattern} documented across the
framework (CLAUDE.md \S1.7; Paper~34
\S\ref{sec:conv_w1e_cross_domain_wall}). The five previous
instances are all in physics observables: Sprint H1 Yukawa (no
autonomous $Y_F$ selection); LS-8a $Z_2/\delta m$ (no autonomous
renormalization counterterms despite faithful reproduction of
two-loop UV-divergent integrand); HF-3 recoil (no cross-register
$V_{ne}(\hat r_e, \hat R_n)$ coordinate operator); HF-4 Zemach
(no magnetization-density operator at the proton coordinate);
HF-5 multi-loop $a_e$ (same renormalization wall as LS-8a, vertex
sector). W1e adds a chemistry-domain instance: the framework
constructs the correct bonding orbital at the FCI level (W1d
closure in F3) but does not autonomously generate the energetic
preference for that orbital at the correct $R_\text{eq}$.
\textbf{The structural-skeleton scope statement extends from
``five physics observables'' to ``five physics + one chemistry
observable'' --- the framework determines the structural skeleton
across both domains but does not autonomously generate calibration
data in either.}

\paragraph{Source documentation (W1e Schmidt + core correlation
closeout).}
Schmidt orthogonalization Day-1 diagnostic STOP:
\texttt{debug/sprint\_w1e\_schmidt\_diagnostic\_memo.md},
\texttt{debug/data/sprint\_w1e\_schmidt\_diagnostic.json}.
[Ne] core-correlation order-of-magnitude estimate DEFER:
\texttt{debug/sprint\_w1e\_core\_correlation\_estimate\_memo.md},
\texttt{debug/data/sprint\_w1e\_core\_correlation\_estimate.json}.

\subsubsection{Structural classification of W1e and the
explicit-core HF decision-gate diagnostic (2026-05-25)}
\label{sec:w1e_classification_and_diagnostic}

After the sprint-scale candidate-space closure
(\S\ref{sec:w1e_schmidt_core_correlation}), a diagnostic-only
structural reading sprint placed W1e within the cross-domain
multi-focal-composition wall pattern and named a load-bearing
decision-gate diagnostic.

\textbf{Structural reading: W1e is H1-class, not LS-8a-class.}
Three independent comparisons place W1e in the same structural
sub-class as Sprint H1's Yukawa-undetermined verdict (admits
structure, lacks autonomous selection) rather than the LS-8a /
HF-5 sub-class (requires external counterterms parameterizing UV
physics): (i) the framework CAN build the right structure --- W1d
closure in F3 constructed the bonding orbital at the FCI level
(naturals $[1.9991, 0.0007]$, four orders advance over W1c+mz
alone); (ii) the closure mechanism is ``external preference data''
(supplying the Pauli pressure to select the bonding orbital at the
correct $R_\text{eq}$), not ``external counterterms'' (which have
the open-ended Landau-pole character of LS-8a / HF-5); (iii) the
calibration data is framework-internally computable in principle
(existing \texttt{geovac/neon\_core.py} infrastructure can be
inverted from screening-potential mode to explicit-orbital mode;
\texttt{geovac/multi\_zeta\_orbitals.py} provides the basis
infrastructure). The chemistry-side distinguishing feature is that
the relevant external input (atomic-structure data) is
\emph{closer at hand} than UV-completion physics --- multi-week to
multi-month architectural commitments rather than a fundamental
UV completion outside the framework.

\textbf{Closure-prospect classification: (a-with-caveat)
sprint-scale exhausted but multi-month tractable for components,
with a calibration-external residual whose magnitude depends on
the Reading A vs Reading B verdict.} The three sub-sub-mechanisms
identified at F4-F6 (basis-truncation $\sim 10.2\%$ / Hartree-pressure
$\sim 25.7\%$ / deep-correlation-residual $65$--$90\%$) stack to
$\sim 35$--$45\%$ combined ceiling at sprint scale; the question is
whether multi-month implementation pushes this to $\sim 85$--$95\%$
(Reading A: active-space deficit, well-defined CCSD-class closure)
or to $\sim 30$--$50\%$ leaving a structurally bounded residual
(Reading B: architectural-scope boundary analogous to H1's
Yukawa-undetermined verdict).

\textbf{Explicit-core 12-electron HF on NaH (~2--4 weeks) is the
named load-bearing diagnostic.} Implement HF on the full 12-electron
NaH problem in the framework's existing basis machinery (single-zeta
CR67 + multi-zeta for valence, $\sim 120$-orbital matrix size at
$n_{\max} = 2$). Run 8-point PES sweep. Decision-gate: if HF gives
internal minimum at $R_\text{eq} \in [3.0, 4.5]$~bohr with
$D_e \in [0.04, 0.12]$~Ha $\rightarrow$ \textbf{Reading A confirmed}
(W1e is active-space deficit; multi-month CCSD-class post-HF closure
path is well-defined). If HF still gives monotone-descending PES
$\rightarrow$ \textbf{Reading B confirmed} (W1e is structurally tied
to the composed-architecture choice; chemistry-domain
architectural-scope boundary established, structurally analogous to
H1's Yukawa-undetermined verdict in the physics domain).

\textbf{Three-class partition of the framework's external-input
space (Sprint Chemistry Positioning Track 2, 2026-05-25).}
The cross-domain wall pattern at the boundary (6 instances, 5
physics + 1 chemistry) is one structural reading. What sits
\emph{outside} the boundary refines into three categorically
distinct classes: \textbf{(1) Calibration parameter values}
(Yukawas, $Z_2$, $\delta m$; UV-physics or AC-inner-factor tied;
structurally bounded; ``second packing axiom'' is the speculative
path forward) --- members H1, LS-8a, HF-5. \textbf{(2) Multi-focal
composition} (spatial $\times$ spatial composition rules; architectural
extension well-defined within existing computational primitives) ---
members HF-3, HF-4 (operator-level largely closed by W1b extension).
\textbf{(3) Multi-determinant correlation} (N-representable
density-matrix manifold beyond the zero-entropy class; chemistry-methods
on the correlation submanifold; ambiguous between Reading A
framework-internal-extendable and Reading B
architecturally-scoped) --- member W1e. The diagnostic-before-engineering
implication: don't propose Class-2 closure mechanisms (single-particle PK,
mean-field Hartree, magnetization extensions) for Class-3 walls --- they
project onto the zero-entropy density-matrix submanifold and saturate
below the wall, as F4-F6 + Schmidt + core correlation empirically
demonstrated.

\paragraph{Source documentation (structural classification + diagnostic).}
W1e structural reading + closure-prospect classification:
\texttt{debug/sprint\_chemistry\_position\_w1e\_structural\_memo.md}.
Modular $\rightarrow$ chemistry clean-negative + three-class partition:
\texttt{debug/sprint\_chemistry\_modular\_connection\_memo.md}.
Memory: \texttt{memory/external\_input\_three\_class\_partition.md}.

%% ========================================================================
\section{Discussion}
\label{sec:discussion}
%% ========================================================================

The composed framework's PK pseudopotential was a necessary compromise
for \emph{classical} solvers that could not handle the full
inter-group Coulomb interaction.  In second quantization on a quantum
computer, antisymmetry is enforced by the fermionic algebra, not by
coordinate-space projectors.  The occupation number representation
naturally prevents double-occupancy without PK.  This suggests that
PK's \emph{antisymmetry enforcement} role is a classical solver
artifact.  However, PK also approximates the physical core-valence
Coulomb interaction energy---this role must be \emph{replaced} by
explicit inter-block integrals, not simply removed.

Track~CB's negative result confirms this distinction.  The failure
was not because PK was needed for antisymmetry, but because the
replacement was incomplete: cross-block $V_{ee}$ was added without
the balancing cross-center $V_{ne}$.  The 10.9\% error of the fully
decoupled Hamiltonian (better than PK's 15.0\%) reinforces the
interpretation that PK overcorrects the Coulomb interaction, but the
remaining 10.9\% error shows that inter-block coupling \emph{is}
needed for quantitative accuracy.

The four-electron FCI results (Sec.~\ref{sec:results}) confirm the
diagnosis from Track~CB.  PK overcorrects (15.0\% error, unbound).
Removing all inter-block coupling undercorrects (10.9\%, unbound).
Adding only cross-block $V_{ee}$ without $V_{ne}$ is energetically
unbalanced (29.0\%, worst of all).  Only the balanced combination of
all four interaction types---within-block $h_1$ and $V_{ee}$,
cross-block $V_{ee}$, and cross-center $V_{ne}$---produces a bound
molecule with the correct qualitative PES shape.  This suggests that
the physical content of PK is not pseudopotential enforcement of
orthogonality, but rather an incomplete approximation to the sum of
cross-block $V_{ee}$ and cross-center $V_{ne}$.

A practical concern: Track~CB found that adding cross-block ERIs
increases the 1-norm from 37.3 to 85.7~Ha.  Even if PK is
eliminated, the net quantum resource cost could increase if the
balanced coupled Hamiltonian has higher 1-norm than the composed
Hamiltonian with PK classically partitioned (Track~BF: electronic-only
1-norm 33.3~Ha).  The sparsity advantage in Pauli count may be
partially offset by a 1-norm disadvantage.  This tradeoff must be
quantified for the balanced Hamiltonian.

The mathematical infrastructure for computing the missing integrals
exists in the Coulomb Sturmian literature and connects directly to
GeoVac's $S^3$ foundation.  The open question is convergence: does
the Shibuya--Wulfman expansion converge fast enough at molecular bond
lengths to preserve the structural sparsity that is GeoVac's primary
computational advantage?

If the answer is yes, the composed framework would be superseded by
coupled composition: same block structure, same Gaunt selection
rules, same $O(Q^{2.5})$ scaling class, but without PK and with
significantly better accuracy.  Paper~17's composed architecture
would then move to the archive, replaced by the coupled composition
as the production molecular Hamiltonian builder.

If the answer is no, the composed framework with PK remains the
production architecture, and the accuracy ceiling is accepted as the
structural cost of maintaining block-diagonal sparsity.  In either
case, this paper documents the mathematical connections and the
concrete data (Track~CB census, FCI benchmarks) that define the
research frontier.

The $n_{\max} = 3$ results with analytical integrals definitively
resolve the key open questions from Phase~3.  The energy convergence
is excellent: 1.8\% $\to$ 0.20\% ($9\times$ improvement) as the
basis grows from $Q = 30$ to $Q = 84$, with the variational bound
respected.  This rules out the PK-like divergence scenario (where
higher angular momentum channels systematically overcounted
correlation).

The $R_{\rm eq}$ drift, however, is confirmed as structural to the
block decomposition, not numerical.  Analytical integrals (machine
precision via incomplete gamma functions) changed $R_{\rm eq}$ by
only 0.007~bohr at $n_{\max} = 3$ (3.287~$\to$~3.280~bohr), while
the drift from $n_{\max} = 2$ to 3 is $+0.053$~bohr.  This drift
is $3\times$ smaller than PK's $+0.15$--$0.22$~bohr/$l_{\max}$ but
in the same outward direction.  The root cause is the block
decomposition itself: each block uses a different effective nuclear
charge ($Z_{\mathrm{eff}}$), and orbitals on different blocks do not
span the same Hilbert space.  The cross-center $V_{ne}$ partially
bridges this gap but cannot fully compensate for the incomplete
orbital overlap between blocks at finite basis.  The balanced
approach is optimal for single-point quantum simulation at fixed
geometries (0.20\% energy error), but not for PES-based equilibrium
geometry determination.  The analytical evaluation of $V_{ne}$ via
incomplete gamma functions adds the cross-center nuclear attraction
to the algebraic registry (Paper~18, Table~I), confirming that the
balanced coupled framework introduces no new exchange constants
relative to the composed architecture.

\emph{Relationship to the LCAO negative result.}  The coupled
composition program reconstructs a multi-center molecular Hamiltonian
from atom-centered orbitals, which is superficially similar to the
LCAO approach abandoned in v0.9.x (Section~3 of the project
guidelines).  The key difference: LCAO used unstructured graph
concatenation with no natural geometry, producing $R$-independent
kinetic energies.  Coupled composition retains the natural geometry
block structure---each electron group inhabits its own optimal
coordinate system---and Gaunt selection rules enforce sparsity within
and across blocks.  The block decomposition provides physical
structure that LCAO lacked.

\emph{Applications paper.}  Paper~20 presents the GeoVac Hamiltonian
library as a practical tool for the quantum computing community, with
benchmark tables comparing GeoVac and Gaussian resource costs at
matched accuracy, FCI PES characterization, and installation
instructions.  The balanced coupled architecture documented in this
paper (Paper~19) provides the PK-free Hamiltonians benchmarked there.


%% ========================================================================
%% BIBLIOGRAPHY
%% ========================================================================

\begin{thebibliography}{30}

\bibitem{loutey_paper7}
J.~Loutey, ``The Dimensionless Vacuum: Recovering the Schr\"odinger
  Equation from Scale-Invariant Graph Topology,'' GeoVac Paper~7 (2025).

\bibitem{loutey_paper11}
J.~Loutey, ``The Molecular Fock Projection: Diatomic Molecules as
  Prolate Spheroidal Lattices,'' GeoVac Paper~11 (2026).

\bibitem{loutey_paper12}
J.~Loutey, ``Algebraic Two-Electron Integrals on the Prolate
  Spheroidal Lattice,'' GeoVac Paper~12 (2026).

\bibitem{loutey_paper14}
J.~Loutey, ``Structurally Sparse Qubit Hamiltonians from Spectral
  Graph Theory,'' GeoVac Paper~14 (2026).

\bibitem{loutey_paper17}
J.~Loutey, ``Composed Natural Geometries for Core-Valence Diatomics:
  LiH and BeH$^+$ from Discrete Graph Topology,'' GeoVac Paper~17 (2026).

\bibitem{fock1935}
V.~Fock, ``Zur Theorie des Wasserstoffatoms,''
  Z.~Phys.\ \textbf{98}, 145 (1935).

\bibitem{shibuya1965}
T.~I.~Shibuya and C.~E.~Wulfman, ``Molecular orbitals in momentum
  space,'' Proc.\ Roy.\ Soc.\ A \textbf{286}, 376 (1965).

\bibitem{koga1987}
T.~Koga and T.~Matsuhashi, ``Sum rules for nuclear attraction
  integrals over hydrogenic orbitals,''
  J.~Chem.\ Phys.\ \textbf{87}, 4696 (1987).

\bibitem{avery2000}
J.~Avery, \emph{Hyperspherical Harmonics and Generalized Sturmians}
  (Kluwer Academic, Dordrecht, 2000).

\bibitem{avery2004}
J.~Avery, ``Many-center Coulomb Sturmians and Shibuya--Wulfman
  integrals,'' Int.\ J.\ Quantum Chem.\ \textbf{100}, 121 (2004).

\bibitem{hoggan2011}
P.~E.~Hoggan, ``How Exponential Type Orbitals Recently Became a Viable
  Basis Set Choice in Molecular Electronic Structure Work,'' in
  \emph{Solving the Schr\"odinger Equation: Has Everything Been
  Tried?}, edited by P.~Popelier (Imperial College Press, 2011),
  Chap.~7.

\bibitem{avery2012}
J.~S.~Avery and J.~E.~Avery, ``Coulomb Sturmians as a basis for
  molecular calculations,'' Mol.\ Phys.\ \textbf{110}, 1593 (2012).

\bibitem{avery2014}
J.~E.~Avery and J.~S.~Avery, ``Molecular integrals for Slater type
  orbitals using Coulomb Sturmians,''
  J.~Math.\ Chem.\ \textbf{52}, 301 (2014).

\bibitem{herbst2018}
M.~F.~Herbst, J.~E.~Avery, and A.~Dreuw, ``Quantum chemistry with
  Coulomb Sturmians: Construction and convergence of Coulomb Sturmian
  basis sets at Hartree--Fock level,'' arXiv:1811.05777 (2018).

\bibitem{latremoliere2019modular}
F.~Latr\'emoli\`ere, ``The Modular Propinquity,''
  Dissertationes Math.\ \textbf{544}, 1 (2019).
  arXiv:1608.04881.

\bibitem{latremoliere2021dual}
F.~Latr\'emoli\`ere, ``The Dual Modular Propinquity and Completeness,''
  J.~Noncommut.\ Geom.\ \textbf{15}, 347 (2021).
  arXiv:1811.04534.

\bibitem{bunge1993}
C.~F.~Bunge, J.~A.~Barrientos, and A.~V.~Bunge, ``Roothaan--Hartree--Fock
  ground-state atomic wave functions: Slater-type orbital expansions and
  expectation values for $Z = 2$--$54$,''
  At.\ Data Nucl.\ Data Tables \textbf{53}, 113 (1993).

\end{thebibliography}

\end{document}
